\documentclass[11pt,english]{article}

\usepackage{float}
\usepackage{xurl}

\usepackage{xcolor}

\usepackage[T1]{fontenc}
\usepackage[final]{microtype}

\usepackage[normalem]{ulem}

\usepackage{amsmath}
\usepackage{amsthm}
\usepackage{nicefrac}
\usepackage{mathpazo}

\newtheorem{theorem}{Theorem}
\newtheorem{corollary}{Corollary}
\newtheorem{lemma}{Lemma}
\newtheorem{assumption}{Assumption}
\newtheorem*{proposition}{Proposition}
\theoremstyle{remark}
\newtheorem{remark}{Remark}

\usepackage{tikz}
\def\checkmark{\tikz\fill[scale=0.4](0,.35) -- (.25,0) -- (1,.7) -- (.25,.15) -- cycle;}

\usepackage{caption}
  figurewithin=none, tablewithin=none, justification=centering,
  singlelinecheck=true}

\raggedbottom

\usepackage{booktabs}
\usepackage{siunitx}

\usepackage[all,cmtip]{xy}

\usepackage{xstring}
\usepackage{ifthen}
\usepackage{catchfile}
\usepackage{threeparttable}

\usepackage{pifont}\newcommand{\xmark}{\text{\ding{55}}}

\usepackage{multicol}

\usepackage[bottom]{footmisc}



\definecolor{blue}{RGB}{0,92,216}
\definecolor{red}{RGB}{213,40,0}
\definecolor{yellow}{RGB}{240,228,66}
\definecolor{green}{RGB}{0,158,115}
\definecolor{grey}{RGB}{180,180,180}
\definecolor{darkgrey}{RGB}{105,105,105}
\definecolor{white}{RGB}{255,255,255}
\definecolor{purple}{RGB}{199,124,255}
\definecolor{orange}{RGB}{197,90,17}

\newcommand{\grey}[1]{\textcolor{grey}{#1}}
\newcommand{\red}[1]{\textcolor{red}{#1}}
\newcommand{\green}[1]{\textcolor{green}{#1}}
\newcommand{\blue}[1]{\textcolor{blue}{#1}}
\newcommand{\purple}[1]{\textcolor{purple}{#1}}
\newcommand{\orange}[1]{\textcolor{orange}{#1}}

\newcommand{\argmax}[0]{\operatorname{argmax}}
\newcommand{\argmin}[0]{\operatorname{argmin}}
\newcommand{\E}{\mathbb{E}}
\newcommand{\R}{\mathbb{R}}
\renewcommand{\d}{\mbox{ d}}
\newcommand{\Var}{\mbox{Var}}

\usepackage[authoryear]{natbib}
\bibliographystyle{aer}

\usepackage{subfiles}

\usepackage{centernot}
\newcommand{\notimplies}{\centernot\implies}

\renewcommand{\lvert}[1]{\left.#1\right|}
\renewcommand{\rvert}[1]{\left|#1\right.}

\newenvironment{refcomment}
  {\begin{quote}\color{blue}\singlespacing}
{\end{quote}}

\usepackage{subcaption} 

\usepackage{colortbl}

\usepackage{relsize}
 

\usepackage[latin9]{inputenc}
\usepackage[english]{babel}

\usepackage{charter}

\usepackage{geometry}
\geometry{verbose,tmargin=1.125in,bmargin=1.125in,lmargin=1.125in,rmargin=1.125in}
\setlength{\textfloatsep}{12pt plus 2.0pt minus 2.0pt}
\emergencystretch=2em

\usepackage{titlesec}
\titleformat*{\section}{\large\bfseries}
\titleformat*{\subsection}{\normalsize\bfseries}
\titleformat*{\subsubsection}{\normalsize\bfseries}
\titlespacing\section{0pt}{8pt plus 4pt minus 2pt}{8pt plus 2pt minus 2pt}
\titlespacing\subsection{0pt}{8pt plus 4pt minus 2pt}{8pt plus 2pt minus 2pt}
\titlespacing\subsubsection{0pt}{8pt plus 4pt minus 2pt}{8pt plus 2pt minus 2pt}

\finalhyphendemerits=200000
\renewcommand\thesection{\Roman{section}}
\renewcommand\thesubsection{\Roman{section}.\Alph{subsection}}
\renewcommand\thesubsubsection{\Roman{section}.\Alph{subsection}.\arabic{subsubsection}}
\hyphenation{Muk-har-lya-mov}
\hyphenation{Shc-her-ba-kov}

\usepackage{setspace}
\onehalfspacing

\usepackage{fancyhdr}
\pagestyle{fancy}
\fancyhead{}
\fancyfoot{}
\fancyfoot[C]{\thepage}
\renewcommand{\headrulewidth}{0pt}

\usepackage{titling}
\setlength{\droptitle}{-2cm}

\usepackage[unicode=true,bookmarks=true,bookmarksnumbered=false,bookmarksopen=false,
    breaklinks=true,pdfborder={0 0 0},pdfborderstyle={},backref=false,colorlinks=false]{hyperref}

\usepackage{ragged2e}
\usepackage{authblk}
\usepackage{lipsum}
\usepackage{makecell}
\usepackage{multirow}
\usepackage{pdflscape}
\usepackage{tabularx}
\usepackage{rotating}
\usepackage[maxfloats=36]{morefloats}
\usepackage{bbm}

\newcommand{\ssize}{5.5cm}
\newcommand{\mssize}{6cm}
\newcommand{\msize}{7.5cm}
\newcommand{\lsize}{9cm}
\newcommand{\vlsize}{11cm}




\renewcommand\Authsep{\hspace{4em}}
\renewcommand\Authands{\hspace{4em}}
\renewcommand\Authfont{\normalsize}
\renewcommand\Affilfont{\itshape}

 
\begin{document}

\renewcommand*{\thefootnote}{\fnsymbol{footnote}}

\author{

  \makebox[\textwidth][c]{\begin{minipage}{0.33\textwidth}\centering
Mark Egan\thanks{Harvard University, Harvard Business School.
Email: megan@hbs.edu.} \\ \vspace{-1em} Harvard University and NBER

    \end{minipage}
    \hfill
    \begin{minipage}{0.33\textwidth}\centering
Gregor Matvos\thanks{Northwestern University, Kellogg School of Management.
Email: gregor.matvos@kellogg.northwestern.edu.} \\ \vspace{-1em} Northwestern University and NBER

    \end{minipage}
    \hfill
    \begin{minipage}{0.33\textwidth}\centering
Amit Seru\thanks{Stanford GSB.
Email: aseru@stanford.edu.} \\ \vspace{-1em} Stanford University and NBER

    \end{minipage}
}

  \vspace{2em}

\makebox[\textwidth][c]{\begin{minipage}{0.33\textwidth}\centering
Lulu Wang\thanks{Northwestern University, Kellogg School of Management.
Email: lulu.wang@kellogg.northwestern.edu.} \\ \vspace{-1em} Northwestern University

    \end{minipage}

    \begin{minipage}{0.33\textwidth}\centering
Vincent Yao\thanks{Georgia State University, J.
Mack Robinson College of Business.
Email: wyao2@gsu.edu.} \\ \vspace{-1em} Georgia State University

    \end{minipage}
} }

\title{Who Pays for Payments?}

\date{\today\thanks{We thank our discussants Matteo Benetton, Chuqing Jin, Marc Rysman, Zhu Wang, and seminar participants at the Bank of Italy, the Hong Kong Institute for Monetary and Financial Research (HKIMFR), Berkeley, the Carey Finance Conference, Emory, the Federal Reserve Board, Imperial, Northwestern, NBER Household Finance, the Stanford Institute for Theoretical Economics Conference, and Toulouse for helpful comments. Mark Egan is an unpaid member of the Fiserv Small Business Index Advisory Council. Vincent Yao is an unpaid member of the Fiserv Small Business Index Advisory Council and a paid consultant. Fiserv had no role in the design of the study, the analysis, the interpretation of the results, the writing of the paper, the decision to publish, and did not review or approve the content of the paper.}}

\maketitle



\pagenumbering{roman} \thispagestyle{empty} 


\begin{abstract}
\begin{singlespace}

We use novel data on the composition and cost of payments across U.S. merchants to quantify consumer redistribution in the payment system.
  Cards charge interchange fees to merchants to fund consumer rewards.
  When merchants raise prices for all consumers in response to these costs, users of low-cost payment methods (e.g., cash and debit) cross-subsidize high-reward credit card users who shop at the same merchant.
  This standard mechanism implicitly assumes that consumers using different payment methods shop at the same merchants and that merchants face similar fees. 
  We show instead that incidence depends on the joint distribution of payment choices across merchants.
  We document two key forces that shape redistribution.
  First, consumer sorting---where consumers who use different payment methods shop at different merchants---limits the exposure of cash and debit users to the effects of high interchange fees.
  Second, interchange fees vary across merchants; where users of different payment methods overlap, such as at large grocery stores, fees are lower due to sector discounts and private negotiations.
  We embed these forces in a sufficient-statistics framework that maps observed heterogeneity directly into redistribution.
  We estimate that interchange fees transfer approximately \$30 billion every year from cash and debit users to credit card users.
  Consumer sorting and merchant fee heterogeneity reduce regressivity by 25\%, but do not eliminate it.
  Finally, we show that both the Durbin Amendment and the rise of premium credit cards have been regressive, highlighting how policy and innovation can reshape the incidence of platform fees.

\end{singlespace}

\thispagestyle{empty} \setcounter{page}{0}\begin{flushleft}

Keywords: Interchange Fees, Redistribution, Durbin Amendment, Price Discrimination, Payment Systems

JEL Classification: E42, D14, L11, L81.\end{flushleft}

\end{abstract}
\clearpage{}

\pagenumbering{arabic} \renewcommand*{\thefootnote}{\arabic{footnote}}
\setcounter{footnote}{0}



\section{Introduction}

Payment systems redistribute across consumers who choose different payment methods but shop at the same stores.
Transfers arise because payment acceptance costs vary significantly by payment method, yet retail prices do not \citep{Stavins2018}.
Since interchange fees---the major component of acceptance costs---flow back to consumers as rewards, a cross-subsidy emerges: all consumers pay higher retail prices, but the users of high-interchange-fee credit cards capture most of the rewards \citep{carlton1994antitrust}.

This standard view implicitly assumes that consumers using different payment methods shop at the same merchants and that merchants face similar acceptance costs. In practice, however, both assumptions fail: consumers sort across merchants based on payment preferences, and merchants face substantial heterogeneity in interchange fees across sectors, sizes, and negotiated contracts. As a result, the incidence of interchange fees cannot be inferred from average fees or representative consumers, but instead depends on the joint distribution of payment choices across merchants.

In recent years, merchants and policymakers have fought against rising interchange fees through both litigation and legislation.
In 2025, a pending settlement against Visa would allow merchants to selectively decline high-fee premium credit cards, such as the Chase Sapphire Reserve, while still accepting lower-fee cards \citep{Andriotis2025}.
The 2022 Credit Card Competition Act sought to lower interchange fees, particularly for small businesses that cannot negotiate for lower fees from the networks \citep{Andriotis2022}.
The U.S. regulatory environment strongly contrasts with the European Union's interchange fee regulation, which has capped interchange fees at less than one-sixth of the highest fees charged in the U.S., highlighting the wide variation in regulatory approaches across jurisdictions.

Despite this policy attention, there is little consensus on who ultimately bears the cost of interchange fees. Existing debates often frame interchange as either a ``tax'' on merchants or a mechanism to fund consumer rewards. Both views overlook that interchange fees primarily redistribute consumption across consumers, and that the magnitude and direction of this redistribution depend on how consumers and merchants are matched in equilibrium.

This lack of consensus reflects a fundamental measurement challenge. 
Redistribution arises only when different consumers use payment methods with different interchange fees at the same merchant \citep{Gans2018}.
Many existing studies rely on payment survey data that lack information on where consumers shop, implicitly assuming common shopping patterns across consumers who use different payments \citep{Felt.etal2023}.

More broadly, the existing literature treats the incidence of interchange fees as a representative-agent problem, abstracting from heterogeneity in both merchant characteristics and consumer sorting. This abstraction is consequential. If consumers using different payment methods sort into different merchants, or if merchants face heterogeneous fees, then average fee differences are no longer sufficient to determine redistribution. 
Instead, incidence depends on the covariance between payment choices and merchant characteristics---a dimension that has been largely unobserved in prior work, and one that is central to determining the direction and magnitude of transfer, and therefore to understanding who ultimately bears the cost of interchange fees.

In this paper, we use novel merchant-level data from Fiserv to quantify how the payment system redistributes among consumers with different payment preferences.
Our central contribution is to show that the incidence of interchange fees is governed by the joint distribution of payment choices and merchant characteristics, rather than by average fee differences across payment methods.

We proceed in three parts.
First, we document striking heterogeneity in both payment mix and interchange fees across merchants.
Second, we show that this heterogeneity reflects systematic consumer sorting across merchants and variation in merchant-level bargaining power and sectoral pricing. 
Third, we develop a sufficient-statistics framework that maps these empirical moments into welfare and redistribution across consumer groups. 
This approach allows us to quantify redistribution without estimating a full demand system, and to isolate the role of sorting and fee heterogeneity in shaping incidence, under empirically plausible pass-through assumptions that we examine in the data and in a structural model.

Our results highlight the magnitude and limits of redistribution through the payment system.
We estimate that interchange fees generate approximately \$30 billion in annual transfers from cash and debit card users to credit card users.
Overall, cash users pay the equivalent of 26\% higher sales taxes than premium credit card users.
Because credit card use increases with income, this represents a \$9.2 billion annual transfer from low- and middle-income households earning less than \$150,000 in annual income to higher-income households.
These transfers are economically significant, comparable in size (but opposite in direction) to major government programs such as SNAP (\$120bn), the Earned Income Tax Credit (\$57bn), and unemployment insurance (\$40bn).
The transfers here are also large relative to other transfers in consumer financial markets, such as the inter-regional transfer due to the GSE's lack of risk-based pricing (\$15 billion), consumer losses from shrouded credit card fees (\$10 billion), and consumer losses from high-fee index funds (\$20 billion) \citep{Hurst.etal2016, agarwal2015regulating, Brown.etal2024}.

To study the effects of interchange fees, we draw on two comprehensive datasets from Fiserv, a large U.S. merchant acquirer.
The first consists of merchant-level settlement data from all Fiserv merchants from 2006 to 2022, covering merchant-level transaction volumes, counts, and fees by card type.
We focus primarily on a 2022 cross-section of approximately 1 million merchants\footnote{The data allow us to study payment volumes at both the outlet level and the corporate level.
We refer to the latter as a merchant.}, which covers around one-fifth of all U.S. card payments.
The second dataset contains transaction data from Fiserv's Clover platform, covering approximately 800,000 merchants between 2019 and 2022.
Clover is the largest U.S. cloud-based point-of-sale (POS) system and processed over \$133 billion in transactions as of 2020.\footnote{Source: \url{https://web.archive.org/web/20230717211716/https://investors.fiserv.com/static-files/820517ed-48a8-4925-aca7-181065168e25}} A unique feature of the Clover data is that it includes both cash and card payments, allowing us to observe the full composition of payments.
Cash transactions are notoriously difficult to capture in large datasets.
We validate the data against public sources, including IRS data on the firm size distribution, BEA data on state-level GDP, industry trade journals, and the Diary of Consumer Payment Choice \citep{Foster.etal2024}.

In the first part of the paper, we document two basic facts about heterogeneity in the cost and composition of payments across merchants.
Our first fact is that payment composition varies substantially across merchants.
Although credit card transactions account for a little more than half of all card transactions, the standard deviation of the share of credit transactions across merchants is approximately 20 percentage points.
The distribution is also bimodal: some stores have almost all credit card shoppers, while others have few.
This reflects consumer sorting, as consumers with different payment preferences tend to shop at a broad range of merchants.
Because redistribution arises only when consumers using different payment methods shop at the same merchants, this variation limits the extent to which credit card interchange fees redistribute consumption from cash and debit card consumers.

The composition of payments across sectors and geographies is consistent with the fact that higher-income individuals are more likely to use credit cards, especially premium credit cards, and less likely to use debit cards.
We first confirm this relationship between income and payment choice using consumer-level survey data from the Diary of Consumer Payment Choice (DCPC) and MRI-Simmons.
Across sectors, travel and retail tend to feature more credit card expenditures, whereas grocery and gas stations tend to feature more debit card and cash transactions.




Our second fact is that card type, merchant sector, and merchant size play fundamental roles in determining merchant-level interchange fees.
While past work has documented that credit card fees are higher than debit card fees, we show that different types of debit and credit cards incur significantly different fees.
For example, debit cards issued by large banks regulated under the Durbin Amendment charge 0.7\%, whereas those issued by smaller, exempt banks charge around 1.1\%.
These high-fee exempt debit cards account for 40\% of debit transactions.
Within credit cards, basic credit cards charge merchants around 1.7\%, while high-fee, high-reward premium credit cards, often used by higher-income consumers, charge around 2.1\%.
Credit card premiumization has accelerated in recent years.
While premium credit cards represented 15\% of credit card volume in 2006, they now account for 60\% as of 2022.
The fee variation across cards plays an increasingly important role in determining merchant-level fees, as cash use---around 11\% of transaction value in the Clover data---has declined.\footnote{Payment methods in the U.S. have shifted gradually over time, with cash use declining and debit card use steadily increasing.
Cash and check use fell from approximately 13\% of consumer payments in 2019 to around 11\% in 2022.
Since 2006, debit cards have gained share relative to credit cards: while credit accounted for roughly 60\% of card payment volume in 2006, it accounted for only 50\% in 2022. }

Even holding card types fixed, interchange fees also vary across merchants.
Merchants in sectors with high credit card use, such as travel and retail, pay high interchange fees, whereas grocery stores and gas stations pay interchange fees around 50 basis points (30\%) lower.
Large merchants also tend to pay lower interchange rates.
For example, large merchants in the grocery and retail sectors, with annual card sales exceeding $\$1$ billion, tend to pay 50 bps (30\%) lower interchange rates on credit and debit cards compared to smaller grocery stores in the same sector.
These lower rates are the consequence of both publicly posted volume discounts and private negotiations and reflect the bargaining power of larger merchants.

We also provide causal evidence that interchange fees have real impacts on prices and consumption.  Using the Durbin Amendment as a natural experiment, we show that lower interchange fees lead to increased sales and lower prices (Appendix \ref{sec:appendix_durbin}).

Evaluating the welfare effects of interchange fees would typically require estimating a full demand system across consumers, merchants, and payment methods, along with modeling strategic pricing by firms. 
This quickly becomes intractable in settings with rich heterogeneity. 
We show that this complexity is unnecessary. 
Under mild conditions on consumer demand and firms' pricing responses to cost shocks, the first-order welfare effects of interchange fees can be summarized by simple moments of the joint distribution of payment shares across merchants and the corresponding fee schedule. 
These statistics remain valid even when consumers' payment and shopping preferences are arbitrarily correlated. 
Importantly, this approach avoids the need to estimate a full demand system or explicitly model strategic interactions among firms, allowing us to map observed heterogeneity directly into redistribution without imposing strong functional-form assumptions on demand.
Under these assumptions, our sufficient statistic approach is a first-order approximation and holds exactly for small fee changes. 
To assess its performance in richer environments, we calibrate a structural model to evaluate how well the sufficient-statistics approximation captures the welfare effects of interchange under alternative assumptions about demand and pricing.
We find that our sufficient-statistics framework closely approximates the full structural model, indicating that the key forces driving redistribution are well captured by the empirical moments we use, and suggesting that our main results are robust.

Having incorporated merchant heterogeneity, consumer sorting, and the structure of interchange fees, we find that the payment system generates large transfers from cash and debit card users to credit card users of around \$30 billion per year.
Our calculations suggest that cash users lose around 96 basis points of purchasing power and that regulated debit card users lose around 47 basis points.
Although much emphasis has been placed on cash users subsidizing credit card rewards, we show that debit card users also subsidize a substantial share of those rewards.
Relative to an average state and local sales tax rate of around 6\%, our calculations suggest that interchange fees are analogous to raising the sales tax rate for cash users by around 16\% and that for regulated debit card users by 8\%.
In contrast, basic and premium credit card users consume around 54 basis points more as a result of high interchange fees, essentially reducing their effective sales taxes by 9\%.

At the same time, we find that consumer sorting matters for the magnitude of this transfer.
Since cross-subsidization occurs at the merchant level, there is no redistribution if cash, debit card, and credit card users shop at different merchants.
In fact, consumer sorting reduces the effective size of the transfer by approximately \$8.6 billion per year.
This effect is largely driven by the tendency of many high-income, premium credit card users to shop at merchants that primarily serve similar customers.
If a merchant primarily receives payments from premium credit cards, there is little to no cross-subsidization.
As a result, consumer sorting insulates cash and debit card users from the full effects of the rise in premium credit card usage.

Variation in interchange fees across merchants, driven by a combination of sector discounts and the bargaining power of large merchants, also shapes the redistributive effects of interchange fees.
Redistribution requires that consumers with different payment methods shop at the same merchants, and that the payment methods incur different interchange fees.
Grocery stores, gas stations, and large retailers exhibit the greatest overlap between consumers using different payment methods, but they also face the lowest levels of interchange fees and the smallest fee differentials across payment methods, due to sectoral discounts and private negotiations. This substantially reduces the extent of redistribution.  
Even though fee heterogeneity suggests that Visa and Mastercard exert market power over merchants, the networks do so in a way that moderates the amount of redistribution in the payment system.
These sector and negotiated discounts lead to a progressive redistribution of approximately \$1.7 billion per year to cash and regulated debit card users, thereby attenuating the overall transfer, highlighting how equilibrium pricing both reflects and mitigates cross-subsidization.

We also use our framework to examine how two recent industry shifts, the Durbin Amendment and the rise in card premiumization, have affected redistribution.
Although cash users benefited from the Durbin Amendment, our analysis shows that credit card users were the primary beneficiaries.
Our results also indicate that the Durbin Amendment was a regressive policy, benefiting wealthier households on net at the expense of lower-income households.
Similarly, we find that the increase in credit card premiumization has been regressive, since higher-income consumers are the most likely to use premium cards. These two examples illustrate how distinguishing between different types of debit and credit cards is essential to understanding the effects of interchange in modern payment systems, and more broadly how policy and innovation can reshape the incidence of platform fees.



\subsection{Related Literature}

Our paper contributes to the literature on two-sided markets by quantifying how payment platform fees redistribute consumption among consumers.
Prior theoretical work has focused on whether high interchange fees encourage efficient payment choices \citep{Rochet.Tirole2011,Wright2012,Edelman.Wright2015}.
These papers typically assume that consumers using different payment methods shop at the same merchants, and that all merchants pay the same interchange fees.
In contrast, we document substantial heterogeneity in both payment composition and interchange fees across merchants and show that this has important implications for redistribution across consumers, implying that incidence cannot be characterized using representative consumers or average fees.

Our framework builds on the literature using sufficient statistics to evaluate the welfare effects of market interventions \citep{Chetty2009, Kleven2021}.
Whereas sufficient statistics are typically used in settings without market power to quantify the social efficiency of government policies, we show that a sufficient-statistics approach can also be applied to quantify redistribution in an industry model with imperfect competition, and in environments with rich heterogeneity across both consumers and firms.

Our analysis also highlights the redistributive potential of payment system innovations.
Existing work has documented the benefits of expanding payment options, such as the gains from debit card adoption in Mexico \citep{Higgins2024} or the consumer value of payment choice \citep{Alvarez.Argente2025}.
Our premium credit card counterfactual provides a complementary perspective by showing that innovation can redistribute consumption when new payment methods charge high fees to merchants to fund consumer benefits.
This is relevant for understanding the effects of new Buy Now, Pay Later products around the world, which charge merchants even higher fees to fund interest-free payments for consumers \citep{Berg.etal2024}.
Our findings help explain why high-income consumers may resist adopting new low-cost payment innovations, such as FedNow or central bank digital currencies \citep{whited2022will}: because they benefit from existing reward structures, they have limited incentives to switch to lower-cost alternatives.

Our findings on redistribution through the payment system contrast significantly with prior work that emphasizes how unsophisticated consumers cross-subsidize sophisticated consumers when both types of consumers pool on the same financial products \citep{Gabaix.Laibson2006, Fisher.etal2024, Agarwal.etal2022}.
In contrast, redistribution arises in our setting when naifs and sophisticates choose different financial products but pool at the same merchants.
This leads to qualitatively different recommendations for reducing redistribution.
Whereas in most traditional settings separating contracts (between consumer types) can reduce redistribution, in our setting, it is precisely this separation across financial instruments (i.e., payment methods) that drives redistribution, shifting attention from financial contract design to the interaction between consumer sorting and merchant pricing.

Finally, our analysis of the Durbin Amendment's redistributive effects builds on \citet{manuszak2017impact} and \citet{Mukharlyamov.Sarin2025}, who document how banks recouped lost interchange revenue through deposit account fees. We use their finding that the decline in interchange affected debit card consumers to form our estimates of how the Durbin Amendment redistributed across consumer groups, linking changes in pricing on the bank side to changes in incidence across consumers.



\section{Institutional Setting and Data\label{sec:institutional-data}}

In this section, we outline the institutional details and describe our data.
We explain how credit and debit card payment systems operate, focusing on key participants, costs, benefits, and transfers.
We then introduce our two primary merchant-level payment datasets from Fiserv, which allow us to observe how the types and costs of payment methods vary across merchants.
These data show which payments merchants receive and what fees they pay.
Together, these data allow us to link merchant-level payment composition and costs to redistribution across consumers.

\subsection{Payment Processing and Interchange Fees}

A typical card payment involves the customer making a purchase, the merchant accepting the payment, and three key financial intermediaries: the merchant acquirer, the issuer, and the card network.
The issuer, effectively the consumer's bank (e.g., Chase), facilitates the transaction from the consumer's perspective by providing the card and processing the payment.
Merchants, in turn, accept card payments by contracting with a merchant acquirer.
This can be a bank (e.g., Wells Fargo) or a fintech firm (e.g., Fiserv, Square), which supplies payment terminals and facilitates settlement and fund transfers.
The card network (e.g., Visa, Mastercard) is an intermediary between the issuer and the acquirer, routing transaction data, collecting and distributing fees, and coordinating settlement.\footnote{American Express and Discover primarily operate as closed-loop (three-party) systems where the network also serves as the issuer, though they have expanded into open-loop arrangements. Our interchange fee analysis focuses on Visa and Mastercard's four-party model.}

Merchants pay transaction fees to accept cards. These fees vary by payment method due to the involvement of different financial intermediaries and the services bundled with each method (e.g., credit card rewards).
The merchant discount fee is split among the intermediaries, with the primary component being the interchange fee.
The interchange fee is remitted to the issuer and is a primary determinant of both the merchant's cost and the consumer's rewards.
As a result, variation in interchange fees is central to understanding cross-subsidization within the payment system.

Figure \ref{fig:payment-flow} illustrates how fees flow among these parties.
When a consumer makes a card payment, the merchant pays a merchant discount fee to its acquirer.
While the acquirer retains a portion of this fee, it passes most of the fee to the issuer as an interchange fee.
Additionally, both the issuer and the acquirer pay network fees to the card network.
The interchange fee is typically the largest component of the merchant discount fee and funds consumer rewards; consequently, it is the focus of much of our analysis.

The typical merchant discount fee for a credit card transaction in the U.S. is 2.25\%, of which the average interchange fee accounts for about 1.90\%; issuers use most of this interchange revenue to fund consumer rewards averaging around 1.40\% \citep{Nilson2023,Drechsler.etal2025}.
In contrast, network fees average only 0.14\% of the transaction value and do not vary with card type or merchant sector.

Our analysis connects the level and heterogeneity of interchange fees across payment methods and merchants to redistribution across consumers.
Interchange fees vary based on the type of card used, the sector and size of the merchant, and ticket size.
Table \ref{tab:visa-interchange-public} presents an excerpt from Visa's 2022 public interchange schedule, illustrating variation along these dimensions.
For example, a premium credit card transaction at a small grocery store incurs an interchange fee of 7 cents plus $2\%$ of the transaction value.
For a $\$100$ transaction, this would amount to $\$2.07$.
The same transaction on a regulated debit card would cost just 27 cents.
If the $\$100$ transaction occurred at a gas station, the cost would be $\$1.10$, while at a large grocery store, it would be $\$1.45$.
While public schedules illustrate potential fee variation, to the best of our knowledge, we are the first to use data on the actual fees merchants pay to quantify how this heterogeneity shapes redistribution.

\subsection{Data}

We use two proprietary datasets from Fiserv, one of the largest U.S. merchant acquirers, to examine how interchange fees vary across merchants and their impact on consumer welfare.
A key feature of these data is that we observe both the composition and price of transactions at the merchant level, which is critical for measuring heterogeneity in payment costs and its implications for redistribution.

\subsubsection{Sources}

Our primary dataset comes from Fiserv's payment settlement records, which cover one-fifth of all U.S. card volume.
Fiserv is a merchant acquirer responsible for processing transactions and calculating the interchange fees owed to issuers, making it an ideal source for studying how the composition and cost of payments vary across merchants.
The data include both in-person and online transactions.

Our data enable us to observe total payment values, payment counts, and interchange fees paid for each card type by each establishment, which we define as the combination of the merchant, sector, and location.
All sales and interchange values are measured net of returns and chargebacks. We observe data over the period 2006-2022. In our cross-sectional analysis, we focus on the most recent cross-section of data (2022), which reflects current pricing structures and consumer behavior.\footnote{We use 2022 as our primary cross-section because it represents the most recent year with complete overlap between our Fiserv settlement data and Clover point-of-sale data. This overlap is essential for observing cash transactions, which are only available in the Clover data.} 
We use merchants' MCC codes to classify establishments into sectors (e.g., grocery stores, gas stations, restaurants).

We study the data at two levels of aggregation. We first aggregate to the establishment-by-sector level to study the geography of payments.
This also allows us to split large multi-sector and multi-location merchants (e.g., those selling gas and retail goods under different MCC codes).
However, because interchange fees are set at the firm level, we aggregate to the firm level for the redistribution analysis, ensuring that fee variation reflects the relevant pricing unit.

Our second dataset comes from Clover, a Fiserv-owned point-of-sale (POS) system, covering approximately 800,000 merchants from 2019 to 2022.
The Clover data uniquely includes cash and check transactions, which are typically unavailable in settlement data.
These additional observations allow us to capture the complete set of transactions at the merchant level, including payment methods that do not incur interchange fees.

Like the settlement data, the Clover dataset contains information on payment values, counts, and interchange fees across card types at the establishment level.
We also observe each merchant's location, sector, and firm size.
While the Clover sample skews toward smaller merchants, it provides valuable insight into cash usage and variation in payment composition that informs redistribution, particularly for merchants where cash remains a meaningful share of transactions.

We supplement our data from Fiserv with survey data on consumers' preferred payment methods from the Federal Reserve Bank of Atlanta's DCPC and the MRI-Simmons Ultimate Survey of Americans (USA).
The Diary of Consumer Payment Choice records the transactions of around 4,000 respondents each year and tracks their payments over a three-day window in October.
Crucially for our study, it provides us with data on consumers' incomes and their preferred payment methods.
The MRI-Simmons USA surveys around $50,000$ respondents each year on their brand preferences, including for financial products.
The data on bank choice and credit card type allow us to group consumers by more granular payment preferences that are relevant for the rewards they receive and the interchange fees they impose on merchants, allowing us to map payment choices to income and reward structures.

\subsubsection{Data Coverage and Representativeness}

The Fiserv data are highly representative of the broader U.S. economy in terms of sectoral coverage, firm size (based on card sales), and geographic distribution.
Figure \ref{fig:data-representative-sector} compares the sector composition of card transactions in the Fiserv data to that reported in the Diary of Consumer Payment Choice.
We observe close alignment across major retail categories, such as restaurants, grocery stores, merchandise retailers, gas stations, and travel.
Figure \ref{fig:data-representative-geo} compares the state-level distribution of card transactions in the Fiserv data with states' shares of GDP.
Most points lie along the 45-degree line, indicating strong geographic representativeness.
Figure \ref{fig:data-representative-size} compares firm size distributions in the Fiserv dataset with those reported by the IRS Statistics of Income, again showing close correspondence across revenue categories.

The Fiserv payment data are also representative of external measures of average transaction size, merchant fees, and the composition of payments.
Figure \ref{fig:fiserv-nilson-representativeness} compares the Fiserv data with aggregate benchmarks from the Nilson Report along four dimensions: transaction size, average fees by card type, the share of card transactions on credit cards, and the share of consumer purchases on cash.
We find a close correspondence along all four dimensions, supporting the external validity of our measures of payment composition and costs.



\section{Variation in Payment Mix and Payment Costs\label{sec:costs}}

Redistribution in the payment system requires consumers who use different payment methods to shop at the same stores.
The magnitude of redistribution at each merchant depends on the differences in the fees for these payment methods.
In this section, we document novel facts about the variation in payment composition and costs across merchants.
We show that payment methods vary significantly in their fees, which generates the potential for redistribution.
However, we also document two forces---consumer sorting and merchant fee heterogeneity---that cause aggregate data to overstate the extent of redistribution in the payment system, by limiting the overlap between consumers using different payment methods and compressing fee differences where such overlap occurs.

\subsection{Variation in Payment Mix\label{subsec:payment-mix}}

One notable feature of the data is that payment composition varies significantly across merchants. Figure \ref{fig:pmt_merchant} shows four dimensions of this heterogeneity.
Figure \ref{fig:cash_share} displays the distribution of the share of cash payments across merchants, focusing on those merchants where cash accounts for at least 2\% of sales.
Observations are at the merchant-by-year level using our Clover data, which includes information on cash payments.
We measure the cash share based on the dollar-weighted share of cash and check payments relative to the total transaction value.
We group cash and check payments into a single category, labeled as ``cash'' because they are distinct from card payments in that they do not incur interchange fees for merchants and do not generate explicit rewards for users.
On average, cash accounts for around 11\% of transaction dollars (dollar-weighted across merchants; Figure \ref{fig:cash_share_ts}).
However, this masks substantial heterogeneity.
For approximately two thirds of merchant-year observations, cash accounts for less than 2\% of transactions.
In contrast, for the one-third of merchants for which cash represents at least 2\% of sales, it accounts for more than 30\% of their transactions---and for merchants in the 90th percentile, over 80\%, highlighting the concentration of cash usage among a subset of merchants.

This heterogeneity shapes redistribution.
For the two thirds of firms where cash use is almost nonexistent, there is little scope for redistribution between cash and card users.
Similarly, at firms where cash is the predominant payment method, redistribution is limited.
Redistribution inherently requires a mixed payment environment, so dispersion in payment composition directly constrains the scope for cross-subsidization.

The remaining 89\% of sales occur via debit and credit card payments.
Here, too, the payment mix varies across merchants.
Figure \ref{fig:credit_share} displays the share of credit card payments relative to total card payments (i.e., debit plus credit) in card settlement data.
Observations are at the merchant-by-month level.
On average, credit card sales account for 53\% of card transactions at the typical merchant in our sample.
Moreover, the share of credit card transactions has remained relatively constant over time (Figure \ref{fig:credit_share_ts}).

The distribution of credit card sales is bimodal, with peaks at approximately 25\% and 70\%.
This bimodality suggests that although credit cards account for 53\% of card transactions on average, merchants tend to fall into two distinct groups: those where credit accounts for around 25\% of card transactions and those where it accounts for about 75\%.
This pattern has implications for cross-subsidization.
Debit cards carry lower interchange fees and rewards relative to credit cards.
Just as cash transactions can subsidize credit card rewards, debit card transactions also potentially play a subsidizing role.
The bimodal nature of the distribution suggests that variation in card payment mix is significant at the merchant level---again limiting the potential for cross-subsidization, as many merchants are dominated by a single payment type.

\subsubsection{Variation Across and Within Sectors}

Part of this heterogeneity in the merchant payment mix stems from sector-specific characteristics.
Figure \ref{fig:pmt_sectors} shows the payment shares, weighted by transaction value, across sectors, for merchants in the Clover data.
Although cash accounts for only 11\% of transactions at the average merchant, it plays a more significant role in sectors such as grocery, restaurants, and gas stations.
For instance, cash accounts for 23\% of sales in the grocery sector, slightly less than the share for credit cards.
By contrast, in the travel sector, credit cards account for 70\% of sales, with debit and cash representing 27\% and 3\%, respectively.
These summary statistics suggest that, in the travel sector, there is limited scope for redistribution between credit and cash users, and any residual redistribution is likely minimal and occurs primarily between credit and debit users, given the limited presence of cash transactions.

We also find substantial dispersion in the composition of payments within sectors.
Figure \ref{fig:sector_densities} plots kernel densities of the share of card transactions on credit cards across merchants in each of the major sectors in the settlement data.
Although grocery stores tend to have a relatively low share of credit card transactions, and travel merchants tend to have a higher share, there remains substantial dispersion within these sectors.
This dispersion is important, as it implies that even within a given sector, consumers with different payment preferences tend to shop at different merchants, reinforcing the role of sorting in shaping redistribution.



\subsubsection{Variation Across Regions and Income Levels}
\label{subsubsec:regional-income-variation}

Beyond sector-level differences, payment methods vary systematically with consumer income, which implies that transfers from cash and debit consumers to credit consumers are regressive.
We combine data from the DCPC and the MRI-Simmons USA to document how payment method preferences vary across the income distribution.
The DCPC provides information on consumer preferences between cash, debit, and credit by income level, while MRI-Simmons USA offers supplementary data on whether debit card users have accounts at small or large banks (relevant for distinguishing exempt versus regulated debit) and whether credit card users hold premium cards (such as Visa Signature or Mastercard World).
Figure \ref{fig:payment_income} illustrates the key pattern: at low levels of income, consumers prefer using debit cards and cash, while at higher levels of income, consumers increasingly use credit cards, and in particular, premium credit cards with higher interchange fees, linking payment choice directly to the distribution of rewards and fees.

Geographic variation in payment mix also reflects these underlying income differences.
Figure \ref{fig:cash_map} displays the county-level share of cash payments.
Light green and yellow regions correspond to higher cash usage, while dark blue areas reflect lower usage.
In some counties (highlighted in yellow), cash accounts for more than 23\% of transactions.
Cash usage is higher in the Midwest and the South, while card usage is more prevalent along the East and West Coasts, as well as in metropolitan areas. In contrast, Figure \ref{fig:credit_map} displays the county-level share of credit card transactions relative to all card payments.
Areas shaded in yellow and light green correspond to higher credit card usage, whereas areas shaded in dark green and blue reflect higher debit card usage.
The results indicate that debit card usage is especially prevalent in the Deep South relative to other U.S. regions.
Conversely, credit card usage is higher in the coastal areas.
These income gradients in payment choices mean that transfers from cash and debit card consumers---driven by interchange fees and rewards---also represent transfers from relatively lower-income consumers to higher-income consumers.

\subsubsection{Decomposing the Variation in Payment Methods}
To better understand what drives the variation in payment methods, we estimate regressions of the following form:
\begin{equation}
Payment\,Share_{jt}=X_{jt}'\beta+\mu_{l(j)}+\mu_t+\mu_{c(j)}+\epsilon_{jt}.

\end{equation}
Observations are at the establishment $(j)$-by-year $(t)$ level.
The dependent variable $Payment\,Share_{jt}$ measures the dollar-share of a given transaction type relative to other payment methods.
We control for merchant characteristics in the vector $X_{jt},$ including average transaction size and total transaction volume.
We also include industry $(l)$, time ($t$), and county $(c)$ fixed effects.

We study two main shares: the share of cash versus all transactions in the Clover data, and the share of credit card transactions versus all transactions in the settlement data. 
We report the corresponding estimates in Table \ref{tab:payment_type}.
Columns (1)--(3) examine cash and check as a share of all payment volumes in the Clover data, while columns (4)--(6) analyze the share of credit card usage within card transactions.
In all columns, the dependent variable is scaled by 100, so coefficients are interpreted in percentage points.

Table \ref{tab:payment_type} shows that consumers are more likely to use cash at merchants with higher sales volumes but lower average transaction values.
In column (1), a 100 log-point increase in total sales is associated with a 1.37 percentage point (12\%) increase in the share of cash transactions.
In contrast, a 100 log-point increase in average transaction value is associated with a 5.55 percentage point (50\%) decline in cash usage.
Cash usage is negatively correlated with e-commerce: a 10 percentage point increase in a firm's e-commerce share is associated with a 0.38 percentage point decline in cash usage.

Consistent with the income gradient documented in Figure \ref{fig:payment_income}, local demographics systematically influence cash usage across merchants. Consumers are more likely to pay in cash in less wealthy, less educated, and less densely populated regions.
Column (2) shows that a one standard deviation increase in median household income is associated with a 0.14 percentage point decline in cash usage.
Similarly, a one standard deviation increase in the share of college-educated households corresponds to a 1.07 percentage point decline in cash usage.

In contrast, consumers are more likely to use credit cards (relative to debit cards) at merchants with lower sales volumes and higher average transaction values.
Column (4) indicates that a 100 log-point increase in transaction value is associated with a 9.11 percentage point (17\%) increase in credit card usage.
Reinforcing the individual-level income patterns we showed earlier, credit card usage is more prevalent in wealthier, older, more educated, and more urban areas, mirroring the distribution of rewards tied to high-fee payment instruments.

Overall, these results suggest that the payment mix varies substantially across merchants, depending on characteristics such as sector and region.
As shown below, this heterogeneity has important implications for both the cost and redistributive aspects of payments, by shaping where cross-subsidization can occur.

\subsection{Variation in Payment Costs}
Payment mix heterogeneity creates substantial variation in merchant transaction costs.
Figure \ref{fig:interchange_dist_card} displays the distribution of average card interchange fees across merchants in our sample.
On average, merchants pay around 1.5\% of sales to accept cards, but this average masks considerable variation.
At the 10th percentile, interchange fees are 0.8\% of sales, while at the 90th percentile, they rise to 2.2\% of sales, implying large differences in effective costs across otherwise similar transactions.

Part of the variation in average transaction costs is driven by the mix of cash versus card payments across firms, as cash transactions incur zero interchange fees, while card interchange fees can exceed 2\%.
However, there is still substantial variation in card interchange fees across merchants, even conditional on payment method.
Figure \ref{fig:interchange_dist_card} displays the average card interchange fee across merchants in our sample with more than \$100,000 in annual sales.
The average fee is 1.5\%, with a standard deviation of 0.5\%.
We next demonstrate that a large part of this variation in fees arises from differences in payment composition, sector-level fee differences, negotiated fees, and ticket size.



\subsubsection{Card Types\label{subsec:card-types}}
Cards vary significantly in their interchange fees.
These differences generate the scope for redistribution, as users of low-fee payment methods cross-subsidize the users of high-fee methods.

We classify card payments into four categories: regulated debit, exempt debit, basic credit, and premium credit.
Figure \ref{fig:price-disc-card} displays average interchange rates for each of these categories.
Regulated debit cards earn around 0.7\% of transaction value in interchange, while exempt debit cards earn higher rates around 1.1\%.
Basic credit cards carry higher interchange rates, around 1.7\%, and premium credit cards have even higher rates, averaging approximately 2.1\%.
Table \ref{tab:icg-summary} shows summary statistics on the distribution of these fees by card type across merchants.

Interchange fees on debit cards depend on whether the issuing bank has more than \$10 billion in assets.
The 2011 Durbin Amendment capped interchange fees for debit cards issued by large banks (regulated debit) but not for those issued by smaller institutions (exempt debit), creating two distinct debit card categories.
Table \ref{tab:public-icg-sector} shows that the regulation reduced debit interchange fees from having a large linear component to a flat rate of around 5 basis points of transaction value plus 22 cents.
Figure \ref{fig:interchange_ts} displays average interchange fees over time, with the drop in debit fees in 2011 reflecting the Durbin Amendment.
Figure \ref{fig:interchange_debit_ts} shows fees separately for regulated and exempt debit cards.

Credit card interchange fees depend on whether the credit card is a premium card targeted at high creditworthiness borrowers.
Premium cards are those in higher interchange tiers, such as Visa Signature, Visa Infinite, Visa Signature Preferred, Mastercard World, and Mastercard World Elite.
These cards offer enhanced rewards and carry higher interchange fees but require consumers to have either higher credit limits or higher levels of spending \citep{Visa2022}.\footnote{Rewards can still vary within card type because the exact structure of rewards is determined by issuers. Visa and Mastercard influence the set of rewards primarily by changing the interchange fees that issuing banks can earn for different card types.}
Basic credit cards include all other credit cards.
Table \ref{tab:public-icg-sector} shows that premium credit cards have publicly reported interchange rates that are higher than those for basic cards.
The increase in average credit card fees shown in Figure \ref{fig:interchange_ts} reflects the growing prevalence of premium credit cards, as the average interchange fees for basic and premium cards have remained relatively stable over this period (Figure \ref{fig:interchange_credit_ts}).
Figure \ref{fig:premium_credit_ts} illustrates the rising dominance of premium cards, which grew from around 15\% of credit card volume in 2006 to 60\% in 2022.\footnote{Historically, card networks such as Visa and Mastercard have required that merchants accepting their cards must accept all cards issued under their respective brands.
In other words, merchants cannot selectively accept certain Visa cards---they must accept all Visa cards or none at all.
This policy has been the subject of ongoing ``honor-all-cards'' litigation \citep{Andriotis2025}.
}

We assess how much differences in payment mix can explain variation in merchant-level interchange fees.
Specifically, we predict each merchant's average interchange fees based on their payment mix across these four types of cards.
Variation in payment mix accounts for approximately 24\% of the overall variation.
The dashed red line in Figure~\ref{fig:interchange_decomp} shows the density of predicted interchange fees based on each merchant's observed payment mix.
This variation in fees can generate variation in retail prices, thereby redistributing consumption across consumers who choose different payment methods.

\subsubsection{Merchant Sectors}

Interchange rates also vary across sectors, with travel and retail typically paying higher fees.
Figure \ref{fig:price-disc-merch} plots average interchange rates in the Fiserv data across five major spending categories: travel, restaurants, merchandise, grocery stores, and gas stations.
These five categories account for approximately 85\% of card sales.
We focus on rates for credit cards and exempt debit cards, as the Durbin Amendment cap limits variation in regulated debit card fees.

Fees are highest in the travel and retail sectors, which generally serve more affluent customers compared to other sectors.
In contrast, grocery stores and gas stations, which serve a broader swath of the income distribution, pay interchange fees that are approximately one-third, or 50 basis points, lower.
Fee variation across sectors reflects Visa's and Mastercard's attempts to balance merchant acceptance versus issuer incentives, trading off participation on both sides of the platform.
Historically, Visa and Mastercard had lower acceptance at supermarkets until they offered discounts starting in the 1990s \citep{Stearns2011}.

Accounting for the merchant sector, in addition to the payment mix, explains roughly 43\% of the overall variation in transaction costs.
The dashed green line in Figure \ref{fig:interchange_decomp} shows the distribution of predicted transaction costs based on each merchant's observed payment mix and sector. 

\subsubsection{Merchant Size and Other Characteristics}

The data show that interchange fees are generally lower for larger firms.
Figure \ref{fig:price-disc-size} plots average interchange rates in the merchandise and grocery store sectors for large firms (with more than \$1 billion in sales) and smaller firms (with less than \$1 billion in sales).
Larger firms pay approximately 50 basis points less in fees on their credit card and exempt debit card transactions.
For credit cards, this pattern aligns with public data from Visa's interchange schedule, shown in Table \ref{tab:public-icg-firm-size}, and likely reflects both publicly posted volume discounts and private negotiations, consistent with greater bargaining power among larger merchants.

A key mechanism underlying these size-based discounts is that large merchants can influence consumer payment behavior by choosing not to accept specific payment methods.
In response, networks offer lower interchange rates to large merchants to keep them on the platform.
Historically, large department stores were slow to accept Visa and Mastercard credit cards because they had their in-house alternatives, such as store credit.
Only after Visa made concessions did JC Penney become the first large department store to accept credit cards, and Sears held out until the mid-1990s \citep{Nocera1995}.

In addition to merchant size, transaction size also plays an important role in determining interchange fees.
Visa's public interchange schedules, as shown in Table \ref{tab:visa-interchange-public}, indicate that interchange fee schedules include a fixed component that does not vary with the transaction size, as well as a variable component that scales linearly with the transaction size.

Accounting for firm characteristics, such as total sales and the average transaction size, in addition to payment mix and sector, explains approximately 64\% of the overall variation in transaction costs.
The dashed orange line in Figure \ref{fig:interchange_decomp} shows the distribution of predicted transaction costs based on each merchant's observed payment mix, sector, and firm characteristics.

\subsection{Real Effects}

In Appendix \ref{sec:appendix_durbin}, we present novel evidence that interchange fees have real effects on merchant sales. To do this, we exploit variation in the effects of the 2011 Durbin Amendment across merchants, which has been exploited in other work such as \citet{Mukharlyamov.Sarin2025}. The Durbin Amendment capped interchange fees only on debit cards issued by large banks, but not on debit cards issued by small banks, nor on credit cards. It also inadvertently increased debit card interchange fees for merchants with small average transaction sizes, while creating significant savings for merchants with large transaction sizes. We implement an instrumental variables design that compares sales growth across merchants with different predicted savings from the Durbin Amendment, using variation in local exposure to large banks to isolate causal effects. We find evidence suggesting that the Durbin Amendment led to a decline in prices and, ultimately, find that every one percentage point in interchange expense saved causes around a 5.6\% increase in card sales. \footnote{Using zip-code level data on retail gasoline prices, \citet{Mukharlyamov.Sarin2025} find that they lack sufficient power to identify the effects of interchange fees on retail prices. In our setting, we are able to use the Fiserv data to show that merchant-level shocks to interchange fees do affect merchant-level sales.}


\subsection{Summarizing the Facts}

Our reduced-form analysis highlights one fact that suggests substantial redistribution between consumers who use different payment methods, and two facts that moderate this redistribution. First, card types vary substantially in their interchange fees. By itself, this raises the potential for substantial redistribution through the payment system. However, redistribution only happens when consumers who choose different payment methods shop at the same stores. Thus, dispersion in payment composition across merchants limits the extent of cross-subsidization. Moreover, at the merchants where consumers with different payment preferences overlap, interchange fees are often lower due to sector discounts and negotiation, further attenuating the magnitude of redistribution.



\section{Methodology: Measuring Interchange Fee Redistribution}
\label{sec:redistribution}

In this section, we describe the development of two related approaches to evaluate the redistributive consequences of interchange fees: a sufficient statistic approach and a structural model.
These approaches allow us to quantify how redistribution changes when accounting for consumer sorting and merchant fee heterogeneity.
We also evaluate different regulatory counterfactuals: European-style regulation and the consequences of the Durbin Amendment.

Each approach rests on different assumptions. The sufficient statistic approach applies broadly across utility functions and market structures, and highlights the first-order determinants of redistribution from interchange fees. A key feature of the sufficient statistic approach is that it imposes minimal structure on the problem. This is a feature because it avoids strong functional-form assumptions; however, without additional structure governing merchant competition, we must make an assumption about how interchange fees pass through to prices. In our baseline implementation of the sufficient statistic approach, we assume full pass-through of interchange fees to prices. We view this as a reasonable approximation for industry-wide cost shocks, but relax it in the structural model. In the structural model, we impose more parametric assumptions on consumer utility functions and the nature of competition between firms. Because the estimates from the two methods are quantitatively similar, we present the more intuitive sufficient statistic approach first.





\subsection{Sufficient Statistic Economic Environment}

We model the behavior of consumers, merchants, and payment networks.
Consumers with fixed payment preferences choose how to allocate spending across merchants subject to a budget constraint.
Merchants set retail prices in response to the composition of payments they receive and the fees on each payment method.
Payment networks set interchange fees and pay consumers rewards.

\subsubsection{Consumers}
The sufficient statistics economic environment places few restrictions on consumer preferences.
Consumers with fixed payment preferences allocate spending across merchants.
Let $k=1,\ldots,K$ denote different payment methods, and let the share of consumers who use payment method $k$ equal $\mu_k$, with $\sum_k \mu_k=1$.
Type-$k$ consumers have income $1 + f_k$, where $f_k$ denotes rewards.
Consumers spend their income on goods from $j=1,\ldots,J$ merchants.
Let $p_j$ denote the price set by merchant $j$ and $q_{jk}$ denote per-capita purchases from merchant $j$ by payment-method-$k$ consumers.
We use stars to denote equilibrium values.

Consumers solve the utility maximization problem

\[q^{*}_{jk} = \arg\max_{q_{jk}\geq 0} U_k\left(\{q_{jk}\}_{j}\right) \text{ s.t. } \sum_j p_j q_{jk} \leq 1 + f_k,\]
where $U_k$ is a strictly increasing, strictly concave, and twice-continuously differentiable utility function.\footnote{ These assumptions rule out perfect substitutes and ensure that demand functions are continuously differentiable almost everywhere.
We assume own-price elasticities dominate cross-price elasticities (Assumption \ref{assump:dominant-diagonal} in Appendix \ref{sec:appendix-theory}) and are not too large relative to the level of merchant fees. } Allowing utility functions to depend on payment preferences enables the model to capture that credit, debit, and cash consumers may shop at different stores.

\subsubsection{Merchants}

Merchants set prices based on their marginal costs to maximize profits.
A merchant's marginal cost depends on expected interchange fees.\footnote{Individual merchants face weak incentives to surcharge by payment method despite different acceptance costs.
Customer backlash creates first-order costs \citep{Stavins2018,Caddy.etal2020}, while a standard envelope theorem argument shows gains from price discrimination are second-order in merchant fees \citep{Wang2025}.}
Let $\tau_{jk}$ denote the interchange fee charged to merchant $j$ for card $k$, expressed as a fraction of the transaction value.  For empirical and expositional convenience, we assume that interchange fees are a function of merchant sector and size. As described in Sections \ref{sec:institutional-data} and \ref{sec:costs}, this assumption is largely consistent with the institutional setting and data.
Let $s_j$ denote the sector of merchant $j$, and let $b_j = 1$ if the merchant has a negotiated discount. The interchange fee charged to merchant $j$ for payment method $k$ is calculated as follows:
\begin{equation}
\tau_{jk}=\overline{\tau}_{k}+\tau_{s_{j}k}+b_{j}\delta_{k}
\end{equation}
where $\overline{\tau}_{k}$ is the baseline fee for payment method $k$, $\tau_{s_{j}k}$ is a sector adjustment, and $\delta_{k}<0$ is a negotiated discount available to large merchants.\footnote{This specification abstracts from the role of ticket size in shaping average transaction fees, as we do not separately model regulatory shocks to the fixed versus linear components of interchange fees.}

Rather than explicitly modeling merchant competition or imposing functional forms on the demand side, we assume that prices are log-linear in weighted interchange fees, such that merchants pass through interchange fees to prices one-for-one, as calculated below:
\begin{equation}
\log p_{j}=\log\overline{p}_{j}+\sum_{k}\mu_{jk}\tau_{jk},
\end{equation}
where $\overline{p}_j$ is a baseline price in the absence of interchange fees, and $\mu_{jk}=\mu_k q_{jk}^{*}/Q_{j}^{*}$ is the share of sales on $k$ where $Q_{j}^{*}=\sum_{k}\mu_k q_{jk}^{*}$ equals total sales at merchant $j$.

We can justify this pricing function in several ways.
First, it is a first-order approximation under CES monopolistic competition.
Second, full pass-through is consistent with models that predict near-complete pass-through for common shocks because all competitors face the same cost increase \citep{Amiti.etal2019}.
Because virtually all major retailers accept cards, interchange fee changes represent sector-wide common cost shocks rather than idiosyncratic shocks.
There are no major retailers unaffected by interchange fees that could constrain pass-through through competitive pressure.
In the restaurant sector, we find causal evidence consistent with this prediction: The Durbin Amendment led to approximately one-for-one pass-through of interchange fees to prices among small firms (Appendix \ref{sec:appendix_durbin}).
We relax this assumption in Section \ref{subsubsec:robustness}.

\subsubsection{Card Networks and Rewards}
We assume that card networks pass-through merchant fees, net of costs, to lump-sum consumer rewards.\footnote{Lump-sum rewards and merchant-specific or sector-specific rewards have identical first-order consumer welfare effects.
By Roy's identity, the welfare impact of reallocating spending across merchants in response to category-based rewards is second-order; the first-order effect is the income change. } 
We interpret rewards broadly to include any subsidy for card use (e.g., cash back, free checking).

Rewards are then determined by:

\begin{align*}
f_{k} & =\frac{\sum_{j}p_{j}q_{jk}^{*}\left(\tau_{jk}-c_{k}\right)}{\sum_{j'}p_{j'}q_{j'k}^{*}}.
\end{align*}
where $c_k$ is the card network's marginal cost associated with payment method $k$. The assumption that issuers fully pass through interchange to rewards is supported by evidence that banks have recouped lost debit interchange via checking fees \citep{Kay.etal2018,Mukharlyamov.Sarin2025} and that credit card issuers pay out a large share of interchange in the form of rewards expense \citep{Drechsler.etal2025,Adams.Bord2020}.\footnote{One concern may be that annual fees on premium cards offset rewards, as described by the influential theory of \citet{Bedre-Defolie.Calvano2013}.
In practice, annual fees primarily screen for high-spending consumers rather than recoup rewards costs: Visa's public interchange schedule distinguishes ``spend-qualified'' rates (Table \ref{tab:visa-interchange-public}), confirming that spending volume directly affects issuer profitability.
Chase lost \$300 million when releasing the Sapphire Reserve because of its generous rewards that more than offset the annual fee \citep{Son2018}.
Further, when Australia cut interchange fees, issuers raised annual fees on rewards cards even though there were less rewards to recoup \citep{RBA2006}.}
Structural models of network competition also predict that networks pass through interchange caps more than one-for-one into reward cuts \citep{Wang2025}.
The assumption can arise in equilibrium if issuers are perfectly competitive, or if they are monopolistically competitive and face logit demand for credit cards with full market coverage.
We relax this assumption in Section \ref{subsubsec:robustness}.

\subsection{Sufficient Statistics Approach}

We develop sufficient statistics to conduct an equilibrium welfare evaluation of interchange fees across consumer groups when payment choices are fixed, but prices and consumption can re-adjust.
The key insight is that first-order consumer welfare effects depend on the variation in the composition and cost of payments across merchants, rather than on demand elasticities. We use this framework to calculate the welfare and redistributive effects of eliminating interchange fees in Section \ref{sec:Results}.

\begin{theorem}[Sufficient Statistic for Consumer Welfare Redistribution]\label{thm:sufficient-stat}
Let $x_l$ be any parameter that shifts $\tau_{jl}$ (e.g. the baseline fee, sector adjustment, or negotiated discount).
The first-order money-metric welfare effect for consumers of type $k$ is

\begin{equation}
\frac{1}{\lambda_k}\frac{dV_k}{dx_l}
= \mu_k^{-1}\!\left(
-\,\mathbb{E}_R\!\left[\mu_{jk}\mu_{jl}\tfrac{\partial\tau_{jl}}{\partial x_l}\right]
+ \mathbb{1}\{k=l\}\,\mathbb{E}_R\!\left[\mu_{jk}\tfrac{\partial\tau_{jk}}{\partial x_k}\right]
\right) + O(\tau_{\max}),
\end{equation}
where $\lambda_k$ is the marginal utility of income, $\tau_{\max} = \max_{j,k} |\tau_{jk}|$ denotes the maximum interchange fee across all merchants and payment methods, and $\mathbb{E}_R\!\left[\mu_{jk}\right]$ denotes expectations weighted by each merchant's share of revenue $\omega_j = R_j / \sum_{j'} R_{j'}$ in the overall economy.
The total effect of a discrete fee component can be calculated by multiplying by $x_l$.
Multiplying by $\mu_k$ yields the aggregate dollar welfare effect for group $k$.

\end{theorem}

\subsubsection{Intuition and proof sketch.}

The theorem states that two channels determine how interchange fees redistribute welfare.\footnote{Our model fixes consumers' payment choices, which allows us to focus on redistribution rather than efficiency effects.
Redistribution can be studied in a largely model-free way by tracking transfers across consumer groups holding behavior fixed, while efficiency analysis requires taking a stand on behavioral responses and therefore demands structural assumptions about demand elasticities, utility functions, and network conduct \citep{Wang2025}.
High interchange fees may induce excessive credit card adoption \citep{Wang2025}, or they may improve efficiency when card acceptance lowers merchant costs \citep{Li.etal2020}.
Our approach enables transparent quantification of redistribution without requiring the stronger assumptions needed for efficiency analysis.}
The first term is the retail-price channel: When interchange fees on payment method $l$ increase, retail prices rise at merchants where method-$l$ users shop.
The size of the mechanical price change, holding fixed the composition of payments, is precisely $\mu_{jl} \frac{\partial \tau_{jl}}{\partial x_l}$.
Consumers using payment method $k$ benefit in proportion to their overlap with method-$l$ users, which then yields the product $\mu_{jk}\mu_{jl} \frac{\partial \tau_{jl}}{\partial x_l}$.
The second term is the direct reward channel: Method-$l$ users receive more in card rewards when their interchange fees rise, affecting only those users.
The sufficient statistic aggregates these price and reward effects across merchants.
Although it is expressed as a mechanical calculation holding payment composition fixed, it captures the full first-order welfare effects even as prices, rewards, and consumption adjust in equilibrium.

The non-trivial part of the theorem is showing that the equilibrium changes in retail prices across all merchants are well approximated by these mechanical effects.
In our setting, merchants' prices depend on the composition of consumers at each merchant.
As prices change, compositions change as well.
While Roy's identity tells us that the effects of consumer reallocation can be ignored when studying the welfare change for a given consumer, it leaves open the possibility that consumer reallocation can change retail prices in a way that differs substantially from the mechanical effect of changing interchange fees on retail prices.
In Appendix~\ref{sec:appendix-theory}, we prove that as long as elasticities satisfy a dominant diagonal condition and are not too large relative to the level of fees, then these feedback effects of prices on the composition of payments are of order $\tau$.
Therefore, when evaluating the redistribution effects of changes in interchange fees $\tau$, the error is second order.

\subsubsection{Relation to alternative approaches.}

Our sufficient-statistic approach yields a simple formula that requires only the payment shares $\mu_k$ and the overlap moments $\mathbb{E}_R[\mu_{jk}\mu_{jl}]$.
In contrast, a fully structural demand estimation requires specifying and estimating utility functions, demand elasticities, and merchant pricing behavior across all consumer types and merchants.
The theorem implies that consumer substitution patterns do not affect first-order redistribution.
To validate this sufficient statistic approach, we next calibrate a structural model that matches key data moments while allowing for imperfect pass-through.
Consumer welfare estimates for the main payment groups change by less than 10\% when allowing imperfect pass-through (Table \ref{tab:compare-ss-model}), confirming that the sufficient statistic formula captures the first-order redistributive effects.

\subsubsection{Sufficient statistic parameters.}

To calculate the sufficient statistics for consumer welfare as per eq. (\ref{eq:sufficient-stat-formula}), all we need are the revenue-weighted first and second-moments of payment shares, $\mathbb{E}_R\!\left[\mu_{jk}\mu_{jl}\right]$ and $\mathbb{E}_R\!\left[\mu_{jk}\right]$, as well as the estimates of the baseline fee $\overline{\tau}_k$, the sector adjustments $\tau_{sk}$, and the negotiated discount $\delta_k$. 

We recover the revenue-weighted first and second-moments with their empirical analogues:
\[
\widehat{\mu}_k = \sum_j \widehat \omega_j\,\mu_{jk}, \qquad
\widehat{\mathbb{E}_R[\mu_{jk}\mu_{jl}]} = \sum_j \widehat \omega_j\,\mu_{jk}\mu_{jl}.
\]
Both sets of moments are straightforward to compute using our merchant-level data.

We recover the fee parameters by regressing firm-level interchange fees for each payment method on sector dummies and a dummy for having more than \$1 billion in sales in the grocery and retail sectors.
We focus on large firms in these sectors, where discounts are most pronounced in the data (Section \ref{sec:costs}).
The regression constant yields the baseline fee $\overline{\tau}_k$; the sector dummy coefficients yield the sector adjustments $\tau_{sk}$; and the bargaining dummy coefficient yields the negotiated discount $\delta_k$.

\subsection{Structural Model}

While the sufficient-statistics approach provides intuitive and transparent estimates of redistribution, it requires us to make assumptions about both the pass-through of interchange fees to retail prices and the pass-through of interchange revenues to consumer rewards.
To assess the robustness of our results to these assumptions, we develop and calibrate a quantitative structural model of the payment system.

The structural model serves two purposes. First, it allows us to assess whether the sufficient statistic's first-order approximation introduces a meaningful error by omitting second-order utility losses from consumer reallocation.
Second, it enables us to evaluate how imperfect pass-through, arising from strategic complementarities in merchant pricing, would affect our distributional conclusions, providing a disciplined way to quantify departures from the baseline assumptions.

We calibrate two versions of the model, where $j^*$ denotes the number of firms per market that internalize their effect on the aggregate price index.
The first version features perfect pass-through ($j^*=50$ firms per market, each with revenue share set near zero so that they behave as price-takers) and validates the accuracy of our first-order approximation.
The second version introduces imperfect pass-through ($j^*=2$ large firms in each sector) and quantifies the sensitivity of our results to this margin.
Both versions are calibrated to match key moments from the Fiserv data on payment composition across merchants and sectors, ensuring that the model reproduces the heterogeneity central to our empirical findings.

\subsubsection{Consumer Preferences}

We model consumer demand using a nested CES structure that allows for rich heterogeneity in shopping patterns across payment methods.
Consumers with different payment methods have different tastes for merchants, and this heterogeneity matches the observed variation in the composition of payments across merchants, linking the model directly to the empirical dispersion documented in the data.

In a representative market $m$, consumers who use payment method $k=1,\dots,K$ have Cobb-Douglas preferences over sectors $s=1,\dots,S$:
\[
\log U_{k}=\sum_{s=1}^{S}\alpha_{ks}\log U_{ks},\quad\sum_{s=1}^{S}\alpha_{ks}=1
\]
Utility in a sector is a CES aggregator over the $J$ firms serving that market:
\[
U_{ks}^{\frac{\sigma-1}{\sigma}}=\sum_{j=1}^{J}a_{jks}^{\frac{1}{\sigma}}q_{jks}^{\frac{\sigma-1}{\sigma}}
\]
where the $a_{jks}$ are taste shifters that allow us to match the distribution of payment shares and firm size. Consumers solve an optimal consumption problem, in which their budget constraint reflects their rewards:
\begin{equation}
\max_{q_{k}}U_{k}\quad\text{s.t. }\sum_{s}\sum_{j}q_{jks}p_{js}\leq1+f_{k} 
\end{equation}

Intuitively, a sector with a high $\alpha_{ks}$ has a high utility for consumers $k$.
For example, premium card consumers may have a higher utility for leisure travel than cash consumers.
Similarly, within a particular sector, consumers' preferences over firms are correlated with their payment methods, which is reflected in $a_{jks}$.
For example, premium card consumers may prefer to shop for groceries at Whole Foods relative to a discount grocer, while cash consumers may prefer the opposite (given prices).
This rich heterogeneity allows firms to have different market shares across consumers who use different payment methods.
The parameter $\sigma$ captures consumers' willingness to substitute across goods (i.e., consumers' price elasticity), so a higher $\sigma$ implies more price elastic consumers. 

The standard solution yields:
\begin{align}
q_{jks} & =\alpha_{ks}\times\left(1+f_{k}\right)\times\frac{a_{jks}p_{js}^{-\sigma}}{P_{ks}^{1-\sigma}} \\ 
P_{ks}^{1-\sigma} & =\sum_{j=1}^{J}a_{jks}p_{js}^{1-\sigma}
\end{align}
The expression is standard and intuitive.
For a given price, goods and sectors with higher $\alpha$ have higher demand.
Higher relative prices lower demand, and higher price elasticity ($\sigma$) results in larger declines in demand given a price increase.

\subsubsection{Merchant Problem}

Merchants set a uniform price for all consumers, regardless of payment method, consistent with the institutional reality that most U.S. merchants do not price discriminate by payment type \citep{Stavins2018}.
A critical feature of our model is that it allows a discrete number of firms in each market, following \citet{Hottman.etal2016}, so that individual merchants internalize their pricing decisions on consumers' demand for goods in that sector.

The departure from the standard monopolistic competition assumption serves two purposes.
First, it allows large merchants to have market power and incomplete pass-through, which is empirically realistic for sectors like grocery and merchandise retail.
Second, it generates strategic complementarities in pricing: when one merchant raises prices in response to higher interchange fees, this partially shifts demand to competitors, dampening their incentive to also raise prices.
By varying the number of large firms per market, we can control the degree of pass-through and thereby assess the sensitivity of our redistribution results to this margin, linking model structure directly to the empirical pass-through question.

Formally, merchants set a uniform price for all of their consumers to maximize profits.
Let the measure of consumers who use payment method $k$ equal $\mu_k$.
For convenience, we normalize marginal costs to 1, as CES models do not separately identify costs and taste shifters and we do not observe cost data:
\[
p_{js}=\argmax_{p}\sum_{k}\mu_k q_{jks}\left(p\left(1-\tau_{jks}\right)-1\right)
\]

Because firms are discrete, they internalize the effect of their own price setting.
Large firms thus exert market power, charge higher markups, and feature incomplete pass-through of interchange fees.
Define $\rho_{jks}$ to be the revenue-share of firm $j$ among all firms in that sector:
\[
\rho_{jks} =\frac{q_{jks}p_{js}}{\sum_{i}q_{iks}p_{is}}
\]

Optimal prices are increasing in the revenue share of the firm, which captures the idea that large firms can exercise market power. 
\begin{align}
  
p_{js} & = \frac{\sum_{k}\mu_{jks}\varepsilon_{jks}}{\sum_{k}\mu_{jks}\left(\varepsilon_{jks}-1\right)\left(1-\tau_{jks}\right)} \\
  \varepsilon_{jks} & = \sigma(1-\rho_{jks}) + \rho_{jks} \\
\mu_{jks} & =\frac{\mu_k q_{jks}}{\sum_{m}\mu_m q_{jms}}
\end{align}

In this model, the pass-through of idiosyncratic cost shocks is decreasing in the revenue share of the firm. When consumers have CES preferences and firms are monopolistically competitive, all merchants fully pass through idiosyncratic cost shocks while ignoring the pricing decisions of their competitors. However, if firms are large, they internalize the effect of their pricing decisions on the aggregate price index; they only partially pass through idiosyncratic cost shocks. Their prices also depend on the prices that other merchants set, i.e., they exhibit strategic complementarity \citep{Amiti.etal2019}. Therefore, although the assumptions we have made about demand rule out perfect pass-through of idiosyncratic cost shocks, strategic complementarities mean that sector-level cost shocks are largely passed through.

\subsubsection{Network Problem}

The network problem in the structural model is the same as that in the sufficient statistics framework. Merchant fees at each firm are based on card type, sector, and firm size, and networks rebate all fees net of costs to the consumer.





\subsubsection{Equilibrium}

An equilibrium is defined by a vector of prices $p_{js}$ for each sector and a vector of consumption choices $q_{jks}$ for consumers of each payment method, such that consumers maximize utility given prices as per Equation \ref{eq:consumption-problem} and merchants maximize profits as per Equation \ref{eq:optimal-price-structural}.

\subsubsection{Calibration}

We next calibrate our model to the data.
We target moments describing the composition of payments across merchants and sectors.
We calibrate most parameters sector-by-sector to allow for the substantial cross-sector heterogeneity documented in Section \ref{sec:costs}.
First, we calibrate the $\alpha_{ks}$ parameters to match the allocation of spending across sectors.
We then parameterize the taste shifters $a_{jks}$ as log-normal, but potentially correlated within firms $j$, to match the revenue-weighted mean and covariance matrices of payment shares across merchants in each sector.
For the grocery and merchandise sectors, we allow for the distributions of $a_{jks}$ to differ for small and large firms (who receive negotiated interchange).
We use the same estimated fee parameters as in the sufficient statistics exercise.
Appendix \ref{sec:appendix-calibrated} provides full calibration details and goodness-of-fit figures, showing that the model closely matches the empirical moments.


\section{Results on Interchange Fee Redistribution}
\label{sec:Results}
We use the sufficient-statistics approach to quantify redistribution under several scenarios.
We first quantify the distributional consequences of interchange fees by measuring redistribution relative to a zero-interchange benchmark, which corresponds to the regulatory environment in the E.U.
We also assess the degree to which the resulting redistribution is progressive or regressive across the income distribution.

We next use a counterfactual to show that the Durbin Amendment, the major regulation of U.S. interchange fees, results in regressive redistribution across consumers, benefiting higher-income households at the expense of lower-income ones.
Given the 2025 North Dakota court ruling (\textit{Corner Post v. Fed}) that declared the implementation of the Durbin Amendment unlawful, this remains a highly relevant policy question.

Finally, we show that the largest losers from the rise of premiumization in the credit card market are debit card users, not cash users, because they frequently shop at the same merchants as premium card users.
This highlights how merchant-level heterogeneity shapes cross-subsidization.

\subsection{Existing Benchmarks}

Before presenting our results, we put them in the context of existing debates.
Public discussion of interchange fees is polarized between two extreme views. 
Merchant groups characterize the fees as a ``hidden tax,'' describing the \$100 billion in annual interchange fees as money siphoned from the economy \citep{MerchantPaymentsCoalition2024}. 
In contrast, bank lobbyists frame interchange solely as funding for consumer rewards, ignoring that the fees can inflate retail prices \citep{BPI2024}. 
Both views miss that interchange fees transfer consumption, rather than imposing pure costs or providing free benefits.

Academic research has attempted to quantify this transfer. Following the logic of \citet{schuh2010gains}, if we focus on credit card versus cash users and assume homogeneous shopping behavior, the calculation is straightforward. U.S. merchants paid approximately \$80 billion in credit card interchange fees in 2022. With credit cards representing roughly half of payment volume, cash users cross-subsidize credit users by approximately \$40 billion annually.

Our analysis yields a more modest estimate of approximately \$30 billion, but more importantly, it reveals a richer picture of who pays whom. Three findings stand out. First, within-card distinctions matter: Regulated and exempt debit cardholders face very different outcomes, as do basic and premium credit cardholders. Second, debit card users---not just cash users---subsidize credit card rewards. Regulated debit cardholders pay almost as much as cash users, a finding obscured by the traditional credit-versus-cash framing. Third, consumer sorting across merchants and merchant fee heterogeneity attenuates redistribution. As \citet{Gans2018} recognized, consumers do not shop randomly; cash and credit card users patronize somewhat different stores. Nevertheless, because we have access to previously unavailable merchant-level data on payment volumes, we can empirically show that sorting reduces but does not eliminate the regressive transfer.



\subsection{The Redistributive Effects of Interchange Fees\label{subsec:eliminate-interchange}}

We study how interchange fees affect retail prices and rewards, and thus redistribution. 
We apply our framework to the 2022 cross-sectional data on sales and fees by card type.
In our main exercise, we measure redistribution relative to a zero-interchange benchmark.
We choose this as a benchmark for two reasons.
First, it corresponds to the regulatory environment in the European Union, which caps most interchange fees near zero, making this counterfactual directly policy-relevant for comparing U.S. and E.U. payment systems.
Second, zero interchange is still consistent with positive merchant fees, as merchant acquirers can still charge a separate margin to cover their own costs. We calculate the welfare effects of eliminating interchange fees using our sufficient statistics framework using eq. (\ref{eq:sufficient-stat-formula}).








\subsubsection{Redistribution Across Payment Types}

The first contribution of our framework is that it allows us to quantify how interchange fees transfer across consumer groups and identify which groups bear these costs.

\paragraph{Cash vs. credit redistribution.}

Table \ref{tab:redist} presents our estimates of redistribution across consumer groups.
We compute the change in consumer surplus arising from interchange fees, accounting for consumer sorting and merchant fee heterogeneity. A positive (negative) value indicates that a consumer group benefits from (loses due to) interchange fees, receiving a positive (negative) net transfer.
Effects are measured both as a percentage of each group's baseline expenditures and in dollar terms when scaled to approximately \$12 trillion in consumer retail purchases \citep{Nilson2023a}.\footnote{Consumer purchases are a subset of PCE since they exclude certain services such as imputed rent and investment management fees that are not paid using cards or cash.}

Cash users lose from interchange fees because these fees raise retail prices at merchants where they shop.
They do not benefit from rewards, since they do not receive them.
On average, cash-using consumers are worse off by about 1 percentage point of consumption.
Across all cash users, this amounts to approximately \$13.9 billion in lost annual consumption.
It is helpful to put these changes in consumption in perspective.
One natural benchmark is the sales tax, which serves as a tax on consumption.
As of 2022, the average state sales tax rate in the U.S. was 6\%.\footnote{\url{https://taxfoundation.org/data/all/state/sales-tax-revenue-reliance-breadth/}} The cross-subsidization between cash and credit card users is equivalent to cash users effectively paying a 26\% higher sales tax than premium credit card users.

\paragraph{Regulated debit card users subsidize credit card rewards.}

Even though much of the literature focuses on transfers between card and cash users, regulated debit card users also lose substantially from interchange fees.
While these debit card holders can earn rewards, these are very small, especially compared to those of credit cards.
The higher retail prices they face far exceed the modest rewards they receive. For this group, average consumer surplus decreases by about 0.4 percentage points of consumption, corresponding to approximately \$13.1 billion in lost annual consumption.
These findings indicate that cash users are not the only ones who subsidize credit card rewards.
Regulated debit card users, who face lower fees and receive lower rewards, collectively contribute a similar order of magnitude in subsidy as cash users.

Whereas cash and regulated debit card users lose from interchange fees, credit card users (basic and premium) gain.
For this group, the rewards from interchange fees outweigh the costs from higher retail prices.
On net, basic and premium credit card users gain 48 and 59 bps of consumption, respectively, amounting to \$9.9 and \$18.2 billion, respectively.
In other words, the current system of interchange fees generates about \$30 billion in transfers from cash and regulated debit card users to credit card users.

\paragraph{Who pays for rewards?} 
We can also quantify who receives and who pays for rewards in Figure \ref{fig:scatter_interchange}.
We calculate how much each payment type contributes to the change in total card rewards when interchange fees increase from zero.
These contributions reflect both card usage and reward levels, expressed relative to the overall change in retail prices.
For example, premium credit card users receive 43\% of rewards but pay only 30\% of total interchange fees.
This 13 percentage point gap represents a transfer to these consumers when they shop at the same merchants with consumers who use other payment methods.
On the other end of the spectrum are cash users, whose share of rewards is zero but who pay for 10\% of interchange fees.
In fact, they implicitly pay for 10\% of card rewards, representing a transfer from cash users to card users.

The figure clearly illustrates why regulated debit card users ``pay'' for rewards of credit card users under the current system. Regulated debit users receive 13\% of rewards but pay 23\% of interchange fees.
Thus, regulated debit cards substantially subsidize the rewards programs associated with more expensive payment methods, such as credit cards.
Exempt debit card users, on the other hand, are almost indifferent to interchange fees as a group.
Their share of rewards is very similar to the implicit cost they pay.
In other words, exempt debit card users pay for their own rewards.

\paragraph{Putting the numbers in context.}

The transfers from interchange fees that we estimate are at least as large as other transfers in credit card financial markets, as well as many government insurance programs.
For credit cards, the consumer losses from shrouded credit card fees are estimated at \$10 billion \citep{agarwal2015regulating,Hurst.etal2016}; transfers arising from credit card interest rates are estimated to redistribute \$18 billion from low credit score consumers to high credit score consumers \citep{Agarwal.etal2022}.
Inter-regional transfers due to the GSEs' lack of risk-based pricing amount to \$15 billion, and consumers' losses from high-fee index funds amount to \$20 billion \citep{agarwal2015regulating,Hurst.etal2016}.
The annual transfers from interchange fees and retail prices are substantial even when measured relative to the outlays of many government social insurance programs, such as SNAP benefits (\$120 billion), EITC (\$57 billion), and regular unemployment insurance (\$40 billion).

\subsubsection{The Effects of Consumer Sorting and Negotiation on Redistribution Estimates}

The second contribution of our framework is that it allows us to quantify the extent to which consumer sorting and merchant fee heterogeneity reduce redistribution. We find that ignoring these two forces inflates the estimated amount of redistribution by roughly one-third. Specifically, we examine the welfare effects of interchange fees relative to a homogeneous benchmark in which: (i) all firms have the same composition of payment methods, and (ii) all firms pay the same interchange fees for a given payment method (i.e., eliminating sectoral variation in interchange fees and negotiated rates).

We decompose the effects of consumer sorting and fee heterogeneity across merchants in Table \ref{tab:redist}. Assuming homogeneous shopping behavior and homogeneous fees across merchants inflates redistribution losses to cash consumers by about one-third and overestimates gains to premium credit card users by about one-half. Formally, under the homogeneity assumption, cash users would lose approximately \$17.7 billion from interchange fees. Similarly, premium credit card users would gain approximately \$27 billion. The homogeneity assumption also significantly overstates the effects on debit card consumers, especially exempt debit users, for whom losses are overstated by a factor of three.

Consumer sorting reduces redistribution across all segments, as shown in the fourth rows of panels (a) and (b) of Table \ref{tab:redist}. Losses to cash and debit card users are reduced, as are gains to all card consumers. Consumer sorting also explains most of the difference between our full model and the model that assumes merchant homogeneity. For example, for premium card consumers, the homogeneous model overestimates gains by \$8.8 billion relative to the heterogeneous benchmark; consumer sorting accounts for \$8.1 billion of this difference. Intuitively, premium cardholders disproportionately shop at premium card-accepting merchants, whereas the reverse is true for cash and debit consumers.

Cash and regulated debit users also benefit from fee heterogeneity. The first column of Table \ref{tab:redist} indicates that consumer sorting and fee heterogeneity reduce the welfare costs associated with interchange fees for cash users by 26 basis points. Consumer sorting accounts for about two-thirds of this effect (18 basis points), as cash users tend to shop at merchants with a higher share of low-cost payment methods. Fee heterogeneity accounts for the remaining one-third of the effect (9 basis points): when cash users shop at the same merchants as consumers who use more expensive payment methods, such as premium credit cards, they tend to do so at merchants---such as large grocery stores---with low negotiated rates. 


\subsubsection{Interchange Fees Result in Regressive Redistribution}

Higher-income consumers use higher-fee payment methods, suggesting that the resulting redistribution is regressive.
Accordingly, we measure the extent of this regressiveness.
We first combine the results on redistribution across payment methods with payment use by income (Figure \ref{fig:payment_income}).
At each income level, we calculate the net effect as a weighted average of the consumption impacts across payment types, using as weights the market shares of each payment method among households at that income level.
We then map these effects to dollar transfers by assuming that consumers of all income levels have the same proportion of purchases to income, and by scaling up total purchases to $\$12$ trillion \citep{Nilson2023a}.

We find that interchange fees generate a large regressive transfer from households with incomes below $\$150,000$ to those with incomes above that level.
In dollar terms, this is a transfer of around \$390 per year to households with incomes above $\$150,000$ per year, funded by households with incomes below that cutoff who lose around \$88 per year.
Cumulatively, high-income households gain around \$9.2 billion from interchange fees, whereas low-income households lose around the same amount.
Our findings help explain why high-income consumers may resist adopting new low-cost payment innovations, such as FedNow or central bank digital currencies: because they benefit from current reward structures, they have weak incentives to switch to low-cost alternatives.

\subsubsection{Robustness}
\label{subsubsec:robustness}

Our benchmark sufficient statistic makes strong assumptions on the pass-through of interchange fees to retail prices and consumer rewards. We also assume that establishment-level prices depend only on the composition of payments at the merchant level. We relax these assumptions below.

\paragraph{Comparison to calibrated structural model.}
Our sufficient statistic assumes that interchange fees are fully passed through to retail prices and that the effects of a discrete change in interchange fees are well approximated by the local effects of small changes in fees.
We systematically evaluate these two potential sources of error with our structural model and find that relaxing these assumptions does not change our finding of substantial redistribution.

As a preliminary step, we first confirm that our calibration successfully replicates the key moments of the empirical data. The first row of Table \ref{tab:compare-ss-model} reproduces our baseline sufficient statistic on the Fiserv data.
The second row computes the sufficient statistic in the synthetic data generated from the calibrated structural model assuming perfect pass-through ($j^*=50$ large firms per market, approximating perfect competition).
The close similarity between the sufficient statistics in the actual data (row 1) and synthetic data (row 2) confirms that the structural model is well calibrated to the moments used in the sufficient statistic.

We next isolate the approximation error from using the sufficient statistic rather than solving the full equilibrium.
Unlike the sufficient statistic, which assumes that the welfare effects of reallocating consumption across merchants are negligible, the structural model allows consumers to fully re-optimize their spending across merchants in response to price changes (row 3).
The similarity in the results in rows 2 and 3 demonstrates that the second-order utility losses from consumer reallocation are negligible.
This validates the envelope theorem intuition underlying our sufficient-statistics approach: Because consumers are optimizing, the marginal benefit of reallocating consumption is zero, so the redistribution effect is well-approximated by the mechanical effect of price and reward changes evaluated at baseline consumption.

Last, we study how modeling imperfect pass-through and market power changes our estimates of redistribution.
In the sufficient statistic, all merchants fully pass through their own costs into higher retail prices.
But when firms have market power, retail prices may reflect both idiosyncratic cost shocks and competitors' pricing behavior \citep{Amiti.etal2019}.
The fourth row of Table \ref{tab:compare-ss-model} shows the results when we calibrate a model with two large firms per sector, in addition to a fringe of monopolistically competitive firms.
Appendix \ref{sec:appendix-calibrated} shows that, in this calibration, the large firms only pass through half of their idiosyncratic cost shocks.
Our evidence from the Durbin Amendment (Appendix \ref{sec:appendix_durbin}) confirms that restaurants (which are typically quite small) appear to pass through interchange fees to prices, consistent with the monopolistically competitive fringe in our structural model.
Taken together, this evidence suggests that the structural model could be seen as a principled way to extrapolate from the pass-through behavior of small merchants to that of larger merchants.

Surprisingly, the structural model finds even greater redistribution from cash and debit card users to credit card users.
Intuitively, strategic complementarity implies that even merchants that do not serve premium credit consumers may raise their prices in response to premium credit interchange fees simply because competitors are raising prices.
This mechanism ``socializes'' the cost of premium credit rewards across all consumers, including those shopping at stores at which premium cards are rarely used.

\paragraph{Imperfect rewards and price pass-through.}
While the structural model captures strategic pricing behavior, it still implies full pass-through of sector-level interchange fees to retail prices.
If consumers can substitute away from retail consumption entirely in response to higher retail prices, then firms may not fully pass through sector-level shocks to interchange into higher retail prices.

We can incorporate imperfect pass-through of retail prices and rewards into our sufficient-statistics framework by scaling the price and reward channels by their respective pass-through rates.
The welfare formula in Theorem \ref{thm:sufficient-stat} has two terms: a price channel (how much consumers pay in higher retail prices) and a reward channel (how much cardholders receive in rewards).
To compute welfare under imperfect pass-through, we scale each term by its respective pass-through rate.
For example, when price pass-through is 70\%, consumers bear only 70\% of the retail price increase implied by interchange fees; when rewards pass-through is 70\%, cardholders receive only 70\% of the rewards implied by full pass-through (see Appendix \ref{sec:appendix-theory} for details).

When both price and rewards pass-through fall to 70\%, redistribution dampens proportionally: The total transfer to credit card users falls to around \$20 billion.
When price pass-through is low (70\%), but rewards pass-through remains full, credit card users benefit more than they would in the baseline: Premium credit users gain \$30.4 billion because they receive full rewards while bearing only partial price increases. 
However, cash and regulated debit card users are still negatively impacted.
When rewards pass-through is low (70\%), but price pass-through remains full, the pattern reverses: Premium credit users gain only \$0.5 billion because they face full price increases while receiving only partial rewards.
Cash and regulated debit users are hurt by more than in the baseline.
These scenarios illustrate that the balance between price and rewards pass-through fundamentally shapes distributional consequences.

\paragraph{Establishment-level shares.}
Our baseline analysis aggregates payment shares to the corporate level, which requires that establishment-level pricing reflects average interchange fees at the corporate level. 
This may overstate redistribution if establishments within the same firm can set retail prices that depend only on the local composition of payments.
The last row of Table \ref{tab:sensitivity-analysis} reports estimates of redistribution if we use the distribution of establishment-level payment shares in the computation of the sufficient statistics.
The estimates are similar to our baseline, suggesting that corporate-level aggregation does not materially affect our conclusions.

\subsection{Effects of the Durbin Amendment}

Our framework also sheds light on the redistributive impact of the Durbin Amendment.
This counterfactual is especially relevant in light of the recent 2025 North Dakota court ruling (\textit{Corner Post v. Fed}) that declared the implementation of the Durbin Amendment unlawful and was temporarily stayed pending an appeal.
The Durbin Amendment was designed to limit interchange fees on debit card transactions.
We compute this counterfactual by replacing all fee components for regulated debit---baseline fees, sector adjustments, and negotiated discounts---with those of exempt debit, then re-normalizing so that the dollar-weighted average sector adjustment remains zero.
Table \ref{tab:durbin} presents estimates of the effects of this regulation.
Regulated debit card users were hurt by the regulation, with their surplus decreasing by 34 basis points of consumption (\$9.6 billion).
Intuitively, prior to Durbin, high debit interchange funded free checking accounts for such consumers, and these perks went away after the regulation \citep{Mukharlyamov.Sarin2025}.
While debit card users were adversely affected, other users benefited from substantially lower retail prices, so the benefit accrued to cash, exempt debit, and credit consumers.
Ultimately, credit consumers' surplus increased by around \$4.3 billion and cash users gained \$1.9 billion, while regulated debit consumers lost \$9.6 billion.

The Durbin Amendment redistributes consumer surplus from middle-income to high-income consumers (Figure \ref{fig:redistribution_durbin}).
Because benefits accrue to both low-income consumers (who use cash) and high-income consumers (who use credit), the net distributional effect is a priori ambiguous.
However, when we apply the payment shares from Figure \ref{fig:payment_income}, we find that low-income consumers were relatively unaffected, middle-income consumers were hurt the most by the Durbin Amendment, and higher-income consumers benefited (Figure \ref{fig:redistribution_durbin}).
This pattern arises because low-income consumers are likely to use both cash and debit cards.
Whereas low-income debit card users were negatively impacted, low-income cash users benefited from lower retail prices.
At the other end of the spectrum, high-income households are most likely to use credit cards, and thus they benefit from the lower retail prices while their rewards are unaffected.
Ultimately, middle-income consumers who primarily use debit cards bore the brunt of the Durbin Amendment.
This implies a cumulative transfer of around \$1 billion from middle-income to high-income consumers.

\subsection{Who Pays for Premiumization?}

Lastly, we examine the redistribution effects of the rise in premium credit cards.
Premium credit cards have recently come into the spotlight due to a proposed settlement that would allow merchants to selectively decline high-fee cards \citep{Andriotis2025}.
As documented in Section \ref{sec:costs} and Figure \ref{fig:premium_credit_ts}, the share of premium credit cards increased from 15\% in 2006 to over 60\% in 2022.
Since the average interchange rate on premium cards is 2.1\%, compared to 1.7\% for basic cards, this shift has raised total credit card interchange fees by roughly 20\%.

Table \ref{tab:premiumization} reports the redistributive effects of the rise in premium credit cards compared to a benchmark where the fees on premium credit cards equal those on basic credit cards.
We compute this counterfactual by replacing all fee components for premium credit---baseline fees, sector adjustments, and negotiated discounts---with those of basic credit, then re-normalizing so that the dollar-weighted average sector adjustment remains zero.
Figure \ref{fig:redistribution_premiumization} displays the corresponding redistributive impact of increased premiumization across household income levels.
Unsurprisingly, premium cardholders benefit from the increase in premiumization: Their surplus increases by 26 basis points of consumption (about \$7.9 billion).
By contrast, cash users, exempt debit users, regulated debit users, and basic cardholders are all hurt by the rise in premiumization.
Notably, the largest costs in dollar terms accrue not to cash users but to debit card users (both regulated and exempt) and basic credit card users.
These groups shop more frequently at the same merchants as premium cardholders, and therefore bear a greater share of the burden from the higher interchange fees associated with premium cards.
Moreover, because all consumers face the same prices regardless of payment method, consumers have incentives to adopt the most expensive card offering the most generous rewards.
Consequently, costly payment methods, such as premium credit cards, tend to crowd out cheaper ones, such as basic cards, helping to explain the rise of premiumization.



\section{Conclusion\label{sec:conclusion}}

Interchange fees shape retail activity and redistribute consumption.
Using merchant-level data, we document substantial heterogeneity in the composition and cost of payments across stores, with critical implications for merchants' costs and redistribution.
By combining the results of a sufficient-statistics framework with a calibrated structural model, we quantify the overall impact of interchange pricing.
Specifically, we find that interchange fees redistribute approximately \$30 billion per year from cash and debit users to credit card users, with consumer sorting attenuating these transfers by about 25 percent.

Our analysis of policy counterfactuals reveals that recent changes in the payment system have had significant regressive impacts.
The Durbin Amendment, despite being designed to help consumers, primarily benefited credit card users rather than the intended debit card users, creating regressive wealth transfers from middle-income to high-income households.
Similarly, the rise of premium credit cards has been regressive, as higher-income consumers disproportionately capture the benefits.
These trends illustrate how network competition through richer rewards funded by higher fees amplifies transfers from low- to high-income households.
More broadly, these results highlight how interchange policies affect not only banks and card networks, but also prices, market competition, and the distribution of economic surplus across consumers.

\clearpage{}
\begin{singlespace}
\begin{flushleft}
{
\bibliography{payment_methods,payments_fiserv_lulu}
}{\par}

\par\end{flushleft}
\end{singlespace}

\clearpage{}


\begin{figure}[H]
\caption{Illustration of Payment Flows\label{fig:payment-flow}}


[Diagram: see PDF version]

\par\textit{Notes:} Figure \ref{fig:payment-flow} displays average fee levels for a credit card transaction in the U.S., showing how fees flow through the payment system. Merchant fees come from \citet{Nilson2023}. Interchange fees and rewards are calculated from regulatory data on large banks \citep{Drechsler.etal2025}. Network fees are obtained from acquirer disclosures and Visa's financial statements \citep{Visa2020,Helcim2021}.

\end{figure}

\clearpage{}
\begin{figure}[H]
  \centering
  \caption{Data Coverage and Representativeness}
  \label{fig:data-representative}

\begin{subfigure}[t]{0.48\textwidth}
    \centering
    \includegraphics[width=\linewidth]{tmp_images/___fiserv_output_graphs_dcpc_settlement_compare_pdf.png}
    \caption{Sector Composition}
    \label{fig:data-representative-sector}
  \end{subfigure}\hfill
  \begin{subfigure}[t]{0.48\textwidth}
    \centering
    \includegraphics[width=\linewidth]{tmp_images/___fiserv_output_graphs_state_scatter_pdf.png}
    \caption{Geographic Composition}
    \label{fig:data-representative-geo}
  \end{subfigure}

\vspace{0.6em}
  \begin{subfigure}[t]{0.48\textwidth}
    \centering
    \includegraphics[width=\linewidth]{tmp_images/___fiserv_output_graphs_firm_size_distn_irs_compare_pdf.png}
    \caption{Firm Size}
    \label{fig:data-representative-size}
  \end{subfigure}

  \par\textit{Notes:} Panel (a) compares sector composition across datasets: DCPC (Atlanta Fed's Diary of Consumer Payment Choice) versus Fiserv settlement data (which covers both large and small merchants). Both show dollar value shares of card transactions by sector. Panel (b) validates geographic coverage by plotting each state's GDP share against its card transaction share in the Fiserv settlement data. Panel (c) compares firm-size distribution of card spending: IRS data (using Statistics on Income for gross receipts) versus Fiserv settlement data. All shares are measured by dollar value of transactions.



\end{figure}\clearpage{}
\begin{figure}[H]
  \centering
  \caption{Heterogeneity in the Payment Mix Across Merchants}
  \label{fig:pmt_merchant}

\begin{subfigure}[t]{0.48\textwidth}
    \centering
    \includegraphics[width=\linewidth]{tmp_images/___fiserv_output_graphs_clover_cash_condl_share_hist_pdf.png}
    \caption{Cash Share}
    \label{fig:cash_share}
  \end{subfigure}\hfill
  \begin{subfigure}[t]{0.48\textwidth}
    \centering
    \includegraphics[width=\linewidth]{tmp_images/___fiserv_output_graphs_clover_cash_share_ts_pdf.png}
    \caption{Cash Share over Time}
    \label{fig:cash_share_ts}
  \end{subfigure}

\vspace{0.6em}
  \begin{subfigure}[t]{0.48\textwidth}
    \centering
    \includegraphics[width=\linewidth]{tmp_images/___fiserv_output_graphs_merchant_cs_credit_share_pdf.png}
    \caption{Credit Card Share of Card Payments}
    \label{fig:credit_share}
  \end{subfigure}\hfill
  \begin{subfigure}[t]{0.48\textwidth}
    \centering
    \includegraphics[width=\linewidth]{tmp_images/___fiserv_output_graphs_credit_share_over_time_annual_sem_pdf.png}
    \caption{Credit Card Share over Time}
    \label{fig:credit_share_ts}
  \end{subfigure}

  \par\textit{Notes:} Panel (a) shows the distribution of cash and check share across merchants in the Clover sample, conditional on at least 2\% cash/check payments; note that two-thirds of merchant-year observations have $\leq$ 2\% cash (nearly all record 0\%), and are excluded from the plot (equally weighted, merchant-year level). Panel (b) displays the dollar-weighted average cash share over time, computed annually. Panel (c) shows the distribution of credit card share of total card transactions in 2022 for merchants with more than \$100,000\ in annual card sales in the broader settlement data. Panel (d) displays the dollar-weighted average credit card share over time, excluding large retailers (grocery/retail >\$1 billion) to avoid entry/exit effects on trends.

\end{figure}\clearpage{}
\begin{figure}[H]
\centering
\caption{Heterogeneity in the Payment Mix Across Sectors} \label{fig:pmt_sectors}

\includegraphics[width=6in]{tmp_images/___fiserv_output_graphs_payment_share_cash_sector_pdf.png}

\par\textit{Notes:} Each bar shows the dollar share of cash, debit, and credit card transactions across major merchant sectors in the Clover merchant data spanning 2019--2022.

\end{figure}

\clearpage{}
\begin{figure}[H]
\centering
\caption{Variation in Composition of Card Payments Within Sectors} \label{fig:sector_densities}

\includegraphics[width=6in]{tmp_images/___fiserv_output_graphs_merchant_cs_credit_share_by_sector_pdf.png}

\par\textit{Notes:} Each density plot shows the distribution of credit card share of total card transactions across merchants in the settlement data in 2022, separately by merchant sector. Sample restricted to merchants with more than \$100,000\ in annual card sales.

\end{figure}

\begin{figure}[H]
  \centering
  \caption{Geographic Variation in the Payment Mix}
  \label{fig:pmt_map}

  \begin{subfigure}[t]{0.9\textwidth}
    \centering
    \includegraphics[width=\linewidth]{tmp_images/map_cash.png}
    \caption{Cash Share}
    \label{fig:cash_map}
  \end{subfigure}

  \begin{subfigure}[t]{0.9\textwidth}
    \centering
    \includegraphics[width=\linewidth]{tmp_images/map_credit.png}
    \caption{Credit Card Share of Card Payments}
    \label{fig:credit_map}
  \end{subfigure}

  \par\textit{Notes:} Panel (a) shows the county-level share of all transactions conducted in cash (equally weighted across Clover merchants in the county). Panel (b) shows the county-level share of card transactions made using credit cards (equally weighted across merchants in the county in the settlement data). Brighter colors indicate higher shares.

\end{figure}

\clearpage{}
\begin{figure}[H]
  \centering
  \caption{Payment Choice and Income}
  \label{fig:payment_income}

    \includegraphics[width=0.9\linewidth]{tmp_images/___output_graphs_payment_income_pdf.png}

  \par\textit{Notes:} Figure \ref{fig:payment_income} shows the share of consumers preferring each payment method---cash, regulated debit, exempt debit, basic credit, and premium credit---by household income bin. Payment preferences are from the DCPC (2022--2023), which records consumers' preferred in-person payment method. The split into regulated versus exempt debit and basic versus premium credit uses card ownership data from the MRI-Simmons survey (2022). Shares are re-weighted within each income bin to match aggregate payment method shares in the settlement data.
\end{figure}

\clearpage{}

\begin{figure}[H]
\centering

\caption{Distribution of Average Transaction Fees (Excluding Cash)} \label{fig:interchange_dist_card}

\includegraphics[width=5in]{tmp_images/___fiserv_output_graphs_merchant_cs_avg_interchange_pdf.png}

\par\textit{Notes:} Figure \ref{fig:interchange_dist_card} shows the distribution of average interchange fees per dollar of card sales across merchants. The sample covers merchants in 2022 with more than \$100,000 in card sales.

\end{figure}

\clearpage{}

\begin{figure}[H]
  \centering
  \caption{Variation in Interchange Fees by Card Type, Merchant Sector, and Size}
  \label{fig:price-disc}

\begin{subfigure}[t]{0.48\textwidth}
    \centering
    \includegraphics[width=\linewidth]{tmp_images/___fiserv_output_graphs_disc_card_type_pdf.png}
    \caption{Interchange Fees by Card Type}
    \label{fig:price-disc-card}
  \end{subfigure}\hfill
  \begin{subfigure}[t]{0.48\textwidth}
    \centering
    \includegraphics[width=\linewidth]{tmp_images/___fiserv_output_graphs_disc_sector_pdf.png}
    \caption{Interchange Fees by Sector}
    \label{fig:price-disc-merch}
  \end{subfigure}

\vspace{0.6em}
  \begin{subfigure}[t]{0.65\textwidth}
    \centering
    \includegraphics[width=\linewidth]{tmp_images/___fiserv_output_graphs_disc_negotiated_pdf.png}
    \caption{Interchange Fees by Merchant Size}
    \label{fig:price-disc-size}
  \end{subfigure}

  \par\textit{Notes:} Panel (a) shows dollar-weighted interchange fees by card type in 2022 for merchants with less than \$1 billion in annual card sales. Regulated debit: cards from banks with >\$10 billion assets. Exempt debit: cards from smaller banks. Premium credit: Visa Signature/Infinite and Mastercard equivalents. Basic credit: all other credit products. Panel (b) shows average interchange fees across merchant sectors for credit and exempt debit cards (excluding firms that receive volume discounts). Panel (c) shows evidence of volume discounts by plotting average interchange fees by merchant size in the grocery and merchandise sectors, including both large merchants (with volume discounts) and smaller merchants (size = total card sales across all merchant establishments).

\end{figure}

\clearpage{}

\begin{figure}[H]
  \centering
  \caption{Interchange Fees by Instrument}
  \label{fig:interchange_dist}

\begin{subfigure}[t]{0.48\textwidth}
    \centering
    \includegraphics[width=\linewidth]{tmp_images/___fiserv_output_graphs_merchant_cs_icg_distn_by_type_pdf.png}
    \caption{Distribution of Fees by Instrument}
    \label{fig:interchange_by_pmt}
  \end{subfigure}\hfill
  \begin{subfigure}[t]{0.48\textwidth}
    \centering
    \includegraphics[width=\linewidth]{tmp_images/___fiserv_output_graphs_debit_credit_icg_rate_annual_sem_pdf.png}
    \caption{Average Debit and Credit Fees over Time}
    \label{fig:interchange_ts}
  \end{subfigure}

\vspace{0.6em}
  \begin{subfigure}[t]{0.48\textwidth}
    \centering
    \includegraphics[width=\linewidth]{tmp_images/___fiserv_output_graphs_debit_icg_annual_sem_pdf.png}
    \caption{Debit Card Fees over Time}
    \label{fig:interchange_debit_ts}
  \end{subfigure}\hfill
  \begin{subfigure}[t]{0.48\textwidth}
    \centering
    \includegraphics[width=\linewidth]{tmp_images/___fiserv_output_graphs_credit_icg_annual_sem_pdf.png}
    \caption{Credit Card Fees over Time}
    \label{fig:interchange_credit_ts}
  \end{subfigure}

  \par\textit{Notes:} Panel (a) shows kernel density plots of average interchange fees by card type in 2022 for merchants with >\$100,000 in sales: regulated debit (leftmost), exempt debit (dotted red), basic credit (dashed green), and premium credit (peaked orange). Panel (b) shows dollar-weighted average interchange fees over time for debit and credit cards combined. Panel (c) shows dollar-weighted debit fees by type over time; the \textit{Regulated} line reflects the Durbin Amendment cap implemented in 2011. Panel (d) shows dollar-weighted credit fees by type over time with premium cards earning higher fees. All time series exclude large retailers (merchandise/retail >\$1 billion) to eliminate entry/exit effects.

\end{figure}

\clearpage{}

\begin{figure}[H]
  \centering
  \caption{Composition of Card Payments}
  \label{fig:card_type_ts}

  \begin{subfigure}[t]{0.7\textwidth}
    \centering
    \includegraphics[width=\linewidth]{tmp_images/___fiserv_output_graphs_debit_reg_share_over_time_annual_sem_pdf.png}
    \caption{Share of Debit that is Regulated}
    \label{fig:regulated_debit_ts}
  \end{subfigure}

  \vspace{0.6em}

  \begin{subfigure}[t]{0.7\textwidth}
    \centering
    \includegraphics[width=\linewidth]{tmp_images/___fiserv_output_graphs_premium_share_over_time_annual_sem_pdf.png}
    \caption{Share of Credit that is Premium}
    \label{fig:premium_credit_ts}
  \end{subfigure}

  \par\textit{Notes:} Data source: Fiserv settlement data. Panel (a) shows the dollar-weighted share of debit card transactions that receive capped interchange rates under the Durbin Amendment over time. Panel (b) shows the dollar-weighted share of credit card transactions that are premium tier, over time. Both time series exclude large retailers (merchandise/retail >\$1 billion) to avoid entry/exit effects on trends.

\end{figure}

\clearpage{}

\begin{figure}[H]
\centering
\caption{Variation in Merchant Interchange Fees Explained by Different Factors}
\label{fig:interchange_decomp}

\includegraphics[width=6in]{tmp_images/___fiserv_output_graphs_merchant_cs_icg_explained_pdf.png}

\par\textit{Notes:} Figure \ref{fig:interchange_decomp} shows the distribution of average interchange fees across merchants (solid blue) and the distribution of predicted fees from progressively richer models. The dashed red line predicts fees from each merchant's card type composition alone (regulated debit, exempt debit, basic credit, premium credit). The dashed green line adds merchant sector. The dashed orange line further adds firm characteristics, including merchant size and average transaction size.

\end{figure}


\begin{figure}[H]
\caption{Who Pays and Receives Rewards}
\label{fig:scatter_interchange}
\centering

\includegraphics[width=5in]{tmp_images/___output_graphs_share_scatter_C_all_with_transfer_pdf.png}
\par\textit{Notes:} Figure \ref{fig:scatter_interchange} shows the share of interchange fees paid and rewards received by consumers across different payment methods (cash, regulated debit, exempt debit, basic credit, premium credit). The figure incorporates merchant heterogeneity in both payment composition and interchange fees and assumes all merchants pass through interchange fees into prices on a one-to-one basis.
\end{figure}

\clearpage{}
\begin{figure}[H]
  \centering
  \caption{Redistribution Across Household Income}
  \label{fig:redistribution}

  \begin{subfigure}[t]{0.9\textwidth}
    \centering
    \includegraphics[width=\linewidth]{tmp_images/___output_graphs_cap_bps_pdf.png}
    \caption{Effects of Current Interchange Fees vs Zero Interchange}
    \label{fig:redistribution_interchange}
  \end{subfigure}

  \begin{subfigure}[t]{0.45\textwidth}
    \centering
    \includegraphics[width=\linewidth]{tmp_images/___output_graphs_durbin_bps_pdf.png}
    \caption{Durbin Amendment}
    \label{fig:redistribution_durbin}
  \end{subfigure}
  \begin{subfigure}[t]{0.45\textwidth}
    \centering
    \includegraphics[width=\linewidth]{tmp_images/___output_graphs_premiumization_bps_pdf.png}
    \caption{Rise of Premium Credit Cards}
    \label{fig:redistribution_premiumization}
  \end{subfigure}

\par\textit{Notes:} Figure \ref{fig:redistribution} illustrates the impact of policy changes on household surplus by income, measured as a percentage of baseline consumption. Panel (a) shows current household surplus relative to a world of zero interchange fees. Panel (b) shows the change in household surplus resulting from the Durbin Amendment. Panel (c) compares current household surplus relative to a world in which premium credit interchange fees equalled that on basic credit cards.

\end{figure}

 
\clearpage{}


\begin{table}[H]
  \centering
  \caption{Summary of Fiserv Data}
  \label{tab:fiserv-summary}

  
    \centering
    \caption{Settlement Data Cross Section from 2022}
    \label{tab:settlement-summary}
    
    {
\def\sym#1{\ifmmode^{#1}\else\(^{#1}\)\fi}
\begin{tabular}{l*{1}{cccccc}}
\hline
                    &           N&        Mean&          SD&         P10&      Median&         P90\\
\hline
Regulated Debit     &   1,410,173&        0.27&        0.22&        0.00&        0.23&        0.57\\
Exempt Debit        &   1,410,173&        0.17&        0.19&        0.00&        0.12&        0.43\\
Basic Credit        &   1,410,173&        0.28&        0.24&        0.03&        0.19&        0.63\\
Premium Credit      &   1,410,173&        0.28&        0.25&        0.00&        0.22&        0.65\\
Annual Card Sales (M)&   1,410,173&        1.33&      182.60&        0.01&        0.09&        0.84\\
Average Interchange (\%)&   1,410,173&        1.49&        0.57&        0.80&        1.47&        2.20\\
\hline
\end{tabular}
}
   

  \vspace{0.8em}

  
    \centering
    \caption{Summary of Clover Data}
    \label{tab:clover-summary}
    
    {
\def\sym#1{\ifmmode^{#1}\else\(^{#1}\)\fi}
\begin{tabular}{l*{1}{cccc}}
\hline\hline
                    &\multicolumn{4}{c}{}                               \\
                    &       Count&        Mean&          SD&      Median\\
\hline
Cash + Check > 2\%  &   1,416,857&        0.36&        0.48&        0.00\\
Conditional Cash + Check Share&     503,876&        0.30&        0.21&        0.25\\
Debit as Share of Card Spend&   1,416,857&        0.52&        0.27&        0.53\\
Annual Card Sales (M)&   1,416,857&        0.27&        0.82&        0.08\\
Annual Sales (M)    &   1,416,857&        0.31&        0.86&        0.09\\
\hline\hline
\end{tabular}
}
   

  \vspace{0.4em}
  \begin{minipage}{\linewidth}
    \textit{Notes:} Panel (a) shows summary statistics from the 2022 cross-section at the
    merchant-sector level. The payment-method variables refer to the share of card sales at the merchant
    on each payment method. Regulated debit cards are issued by large banks and exempt debit cards are issued
    by small banks. Premium credit refers to high-interchange credit cards in the Visa Signature and Visa
    Infinite (and the Mastercard equivalent) interchange tiers. Average interchange is calculated as total
    interchange paid divided by total card sales. Transaction size is the total card dollar volume divided by the
    number of transactions. Panel (b) shows summary statistics of the Clover data from 2019--2022 at the merchant-year
    level. The cash and check share is defined as the dollar value of cash and check transactions divided
    by the total value of transactions for the merchant in that year.
  \end{minipage}
\end{table}

\clearpage

\begin{table}[H]
\caption{Spatial and Firm Determinants of Payment Mix \label{tab:payment_type}}

\begin{centering}
\resizebox{\textwidth}{!}{{
\def\sym#1{\ifmmode^{#1}\else\(^{#1}\)\fi}
\begin{tabular}{l*{6}{c}}
\hline
                    &\multicolumn{1}{c}{Cash}&\multicolumn{1}{c}{Cash}&\multicolumn{1}{c}{Cash}&\multicolumn{1}{c}{Credit}&\multicolumn{1}{c}{Credit}&\multicolumn{1}{c}{Credit}\\
\hline
Share Online        &       -3.76***&       -3.81***&       -3.79***&        2.63***&        3.28***&        3.58***\\
                    &      (0.08)         &      (0.07)         &      (0.07)         &      (0.07)         &      (0.07)         &      (0.07)         \\
[1em]
Log Ticket Size     &       -5.55***&       -4.40***&       -4.38***&        9.11***&        8.29***&        8.20***\\
                    &      (0.01)         &      (0.01)         &      (0.01)         &      (0.01)         &      (0.01)         &      (0.01)         \\
[1em]
Log Sales           &        1.37***&        1.05***&        1.05***&       -1.42***&       -1.16***&       -1.08***\\
                    &      (0.01)         &      (0.01)         &      (0.01)         &      (0.01)         &      (0.01)         &      (0.01)         \\
[1em]
Log Population (Std. Dev.)&                     &       -0.61***&                     &                     &        1.46***&                     \\
                    &                     &      (0.02)         &                     &                     &      (0.02)         &                     \\
[1em]
Median Income (Std. Dev.)&                     &       -0.14***&                     &                     &        1.58***&                     \\
                    &                     &      (0.02)         &                     &                     &      (0.02)         &                     \\
[1em]
Pct. Hispanic (Std. Dev.)&                     &       -0.02         &                     &                     &        2.03***&                     \\
                    &                     &      (0.02)         &                     &                     &      (0.02)         &                     \\
[1em]
Pct. College (Std. Dev.)&                     &       -1.07***&                     &                     &        3.79***&                     \\
                    &                     &      (0.02)         &                     &                     &      (0.02)         &                     \\
[1em]
Pct. 65+ (Std. Dev.)&                     &        0.19***&                     &                     &        3.35***&                     \\
                    &                     &      (0.02)         &                     &                     &      (0.02)         &                     \\
[1em]
Pct. Black (Std. Dev.)&                     &        0.63***&                     &                     &        0.12***&                     \\
                    &                     &      (0.02)         &                     &                     &      (0.02)         &                     \\
\hline
Observations        &   1,416,842         &   1,416,839         &   1,416,793         &   2,172,044         &   2,172,030         &   2,172,038         \\
R2                  &        0.20         &        0.23         &        0.25         &        0.27         &        0.33         &        0.35         \\
DV Mean             &          11         &          11         &          11         &          51         &          51         &          51         \\
Sector FE           &                     &           X         &           X         &                     &           X         &           X         \\
County FE           &                     &                     &           X         &                     &                     &           X         \\
\hline
\end{tabular}
}
 }
\par\end{centering}
\par\textit{Notes:} Dependent variables are firm-level transaction shares. Columns (1)--(3): Cash and check share, estimated on Clover data. Columns (4)--(6): Credit card share of card transactions, estimated on Fiserv settlement data. Independent variables include firm characteristics (sector, size), and local demographics. All regressions are estimated at the establishment level.
\end{table}

\clearpage

\begin{table}[H]
\caption{Summary Statistics on Interchange Rates by Card Type Across Firms
\label{tab:icg-summary}}

\centering

{
\def\sym#1{\ifmmode^{#1}\else\(^{#1}\)\fi}
\begin{tabular}{l*{1}{cccccc}}
\hline\hline
                    &\multicolumn{6}{c}{}                                                         \\
                    &           N&        Mean&          SD&         P10&      Median&         P90\\
\hline
Regulated Debit     &   1,225,787&        0.60&        0.62&        0.09&        0.36&        1.50\\
Exempt Debit        &   1,115,701&        1.43&        0.54&        0.82&        1.43&        2.02\\
Basic Credit        &   1,272,668&        1.64&        0.70&        0.57&        1.80&        2.42\\
Premium Credit      &   1,206,011&        2.34&        0.28&        2.06&        2.36&        2.60\\
\hline\hline
\end{tabular}
}
 
\par\textit{Notes:} Data source: Fiserv settlement data, 2022. Interchange rates calculated as total net interchange divided by total net sales for each card type. Rates are missing if a payment method had no transactions for that merchant. Merchants are equally weighted (not dollar-weighted), which explains the differences in dollar-weighted averages in Figure \ref{fig:price-disc}.
\end{table}

\clearpage

\begin{table}[H]
  \centering
  \caption{Interchange Rates from Visa's 2022 Interchange Publication}
  \label{tab:visa-interchange-public}

  
    \centering
    \caption{By Sector and Card Type}
    \label{tab:public-icg-sector}
    
    \renewcommand{\arraystretch}{1.1}
    \begin{tabularx}{0.95\textwidth}{l*
        {4}{>{\centering\arraybackslash}X}}
      \hline
      Card Category & Gas Station & Supermarket & Merchandise & Restaurant \\
      \hline
      Regulated Debit Card & $0.05\%+0.22$ & $0.05\%+0.22$ & $0.05\%+0.22$ & $0.05\%+0.22$ \\
      Exempt Debit Card    & $0.80\% + \$0.15$ (cap \$0.95) & \$0.30 & $0.80\%+0.15$ & $1.19\%+0.10$ \\
      Basic Credit         & $1.15\% + \$0.25$ (cap \$1.10) & $1.50\%+0.07$ & $1.51\%+0.10$ & $2.10\%$ (floor \$0.04) \\
      Premium Credit       & 1.15\% + \$0.25 (cap \$1.10) & $2.00\%+0.07$ & $2.10\%+0.10$ & $2.60\%$ (floor \$0.04) \\
      \hline
    \end{tabularx}
  

  \vspace{0.8em}

  
    \centering
    \caption{By Firm Size}
    \label{tab:public-icg-firm-size}
    
    \renewcommand{\arraystretch}{1.1}
    \begin{tabular}{l>
        {\centering\arraybackslash}p{1.3in}}
      \hline
      Supermarket Size Category & Premium Credit Rate \\
      \hline
      \makecell[l]{Tier 0 $\left(\geq \$19.6\ \text{Bn}\right)$} & $1.40\% + \$0.05$ \\
      \makecell[l]{Tier 1 $\left(\geq \$10.65\ \text{Bn}\right)$} & $1.55\% + \$0.05$ \\
      \makecell[l]{Tier 2 $\left(\geq \$4.2\ \text{Bn}\right)$} & $1.65\% + \$0.05$ \\
      \makecell[l]{Tier 3 $\left(\geq \$0.95\ \text{Bn}\right)$} & $1.75\% + \$0.05$ \\
      All Other (Smallest) & $2.00\% + \$0.07$ \\
      \hline
    \end{tabular}
  

  \par\textit{Notes:} Data source: Visa's 2022 interchange schedule (publicly available). Fees are shown as variable (percent) plus fixed (\$) components. Example: ``0.05\% + \$0.22'' means 0.05\% of transaction value plus 22 cents. Premium credit = Visa Infinite Spend Qualified Cards. Basic credit = All Other Products category. Regulated debit = debit cards from banks with >\$10 billion assets (subject to Durbin Amendment).
\end{table}

\clearpage

\begin{table}[H]
  \centering
  \caption{Effects of Capping Interchange Fees}
  \label{tab:redist}

  
    \centering
    \caption{As Basis Points of Baseline Consumption}
    \label{tab:redist-se-bps}
    
    \begin{tabular}{l l
                    c c
                    c c
                    c}
      \hline
       &  & \multicolumn{5}{c}{Consumer Group} \\
      \multicolumn{2}{c}{Mechanism} & Cash & Regulated Debit & Exempt Debit & Basic Credit & Premium Credit \\
      \hline
      \multicolumn{2}{l}{Full Model} & -96 (2) & -47 (9) & -4 (10) & 48 (5) & 59 (4) \\
      \multicolumn{2}{l}{Homogeneous} & -123 (4) & -56 (7) & -13 (8) & 49 (5) & 88 (5) \\
      \hline
      \multicolumn{2}{l}{Difference} & 26 (4) & 9 (4) & 9 (5) & -1 (5) & -29 (5) \\
       & Consumer Sorting & 18 (1) & 7 (2) & 13 (2) & -2 (1) & -27 (2) \\
       & Merchant Fee Heterogeneity & 9 (4) & 2 (4) & -4 (5) & 1 (5) & -2 (5) \\
      \hline
    \end{tabular}
  

  \vspace{0.8em}

  
    \centering
    \caption{In Dollar Terms (\$Bn)}
    \label{tab:redist-se-dollars}
    
    \begin{tabular}{l l
                    c c
                    c c
                    c}
      \hline
       &  & \multicolumn{5}{c}{Consumer Group} \\
      \multicolumn{2}{c}{Mechanism} & Cash & Regulated Debit & Exempt Debit & Basic Credit & Premium Credit \\
      \hline
      \multicolumn{2}{l}{Full Model} & -13.9 (0.3) & -13.1 (2.4) & -1.1 (2.9) & 9.9 (1.2) & 18.2 (2.0) \\
      \multicolumn{2}{l}{Homogeneous} & -17.7 (0.6) & -15.7 (2.0) & -3.8 (2.1) & 10.2 (1.2) & 27.0 (1.6) \\
      \hline
      \multicolumn{2}{l}{Difference} & 3.8 (0.6) & 2.6 (1.1) & 2.7 (1.4) & -0.3 (1.1) & -8.8 (1.5) \\
       & Consumer Sorting & 2.6 (0.2) & 2.1 (0.5) & 3.9 (0.5) & -0.4 (0.3) & -8.1 (0.7) \\
       & Merchant Fee Heterogeneity & 1.2 (0.6) & 0.5 (1.0) & -1.2 (1.5) & 0.1 (1.1) & -0.7 (1.7) \\
      \hline
    \end{tabular}
  

  \vspace{0.8em}
 
  \par\textit{Notes:} This table shows the estimated consumer welfare effects of eliminating all interchange fees. Panel (a) presents effects as basis points of baseline consumption. Panel (b) shows effects in dollar terms (normalized to \$12 trillion in total purchases, from the Nilson Report). ``Full Model'' incorporates heterogeneity in both payment composition and costs across merchants. ``Homogeneous'' merchants assumes uniform payment composition. ``Difference'' shows the combined attenuation effect of merchant heterogeneity. ``Consumer Sorting'' and ``Merchant Fee Heterogeneity'' are additive components that together equal the ``Difference'' row. Standard errors for all entries are computed via a parametric bootstrap, described in Appendix \ref{subsec:appendix-bootstrap}.
\end{table}

\clearpage

\begin{table}[H]
  \centering
  \caption{Sensitivity Analysis: Data Aggregation and Pass-Through Rates}
  \label{tab:sensitivity-analysis}
  

  \begin{tabular}{lrrrrr}
\toprule
 & \multicolumn{5}{c}{Consumer Group (\$ billion)} \\
Assumption & Cash & Reg.\ Debit & Exempt Debit & Basic Credit & Premium Credit \\
\midrule
Baseline (Full Pass-Through) & -13.9 & -13.1 & -1.1 & 9.9 & 18.2 \\
70\% Pass-Through (Both) & -9.7 & -9.2 & -0.8 & 6.9 & 12.7 \\
70\% Price Pass-Through & -9.7 & -3.7 & 6.7 & 17.1 & 30.4 \\
70\% Rewards Pass-Through & -13.9 & -18.6 & -8.6 & -0.3 & 0.5 \\
Establishment-Level Data & -13.5 & -13.1 & 0.7 & 8.6 & 17.3 \\
\bottomrule
\end{tabular}

 
  \par\textit{Notes:} Table \ref{tab:sensitivity-analysis} presents sensitivity analysis examining how baseline results vary under alternative assumptions. ``Baseline'' replicates our main estimates with full pass-through. ``70\% Pass-Through (Both)'' reduces both price and rewards pass-through to 70\%. ``70\% Price Pass-Through'' reduces only price pass-through to 70\% while maintaining full rewards pass-through. ``70\% Rewards Pass-Through'' maintains full price pass-through but reduces rewards pass-through to 70\%. ``Establishment-Level Data'' uses establishment-level rather than firm-level aggregation.
\end{table}

\clearpage

\begin{table}[H]
\caption{Comparison of Sufficient Statistic and Structural Approaches\label{tab:compare-ss-model}}


\caption{Effects in Basis Points}

\begin{tabular}[t]{lrrrrr}
\toprule
\multicolumn{1}{c}{ } & \multicolumn{5}{c}{Consumer Group} \\
Scenario & Cash & Regulated Debit & Exempt Debit & Basic Credit & Premium Credit\\
\midrule
Sufficient Statistic (Data) & -96.2 & -46.7 & -3.7 & 48.0 & 59.5\\
Sufficient Statistic (Model) & -97.0 & -46.8 & -4.2 & 47.5 & 59.5\\
Structural, Complete Passthrough & -97.5 & -48.0 & -5.3 & 45.6 & 56.4\\
Structural, Incomplete Passthrough & -99.7 & -51.2 & -8.3 & 46.6 & 59.7\\
\bottomrule
\end{tabular}
 

\vspace{0.5cm}


\caption{Effects in Dollars}

\begin{tabular}[t]{lrrrrr}
\toprule
\multicolumn{1}{c}{ } & \multicolumn{5}{c}{Consumer Group} \\
Scenario & Cash & Regulated Debit & Exempt Debit & Basic Credit & Premium Credit\\
\midrule
Sufficient Statistic (Data) & -13.9 & -13.1 & -1.1 & 9.9 & 18.2\\
Sufficient Statistic (Model) & -13.8 & -13.1 & -1.2 & 9.8 & 18.3\\
Structural, Complete Passthrough & -14.0 & -13.5 & -1.5 & 9.4 & 17.2\\
Structural, Incomplete Passthrough & -14.4 & -14.4 & -2.4 & 9.6 & 18.2\\
\bottomrule
\end{tabular}
 

\par\textit{Notes:} Table \ref{tab:compare-ss-model} compares welfare estimates from the sufficient-statistics approach to the exact equilibrium effects from the structural model. Row 1 reports baseline sufficient-statistic estimates from the Fiserv data. Row 2 applies the sufficient-statistic formula to synthetic data from the calibrated model with perfect pass-through ($j^*=50$ firms per market). Row 3 reports exact equilibrium welfare effects from the structural model with perfect pass-through, allowing consumers to fully re-optimize. Row 4 reports exact equilibrium effects with incomplete pass-through ($j^*=2$ large firms in each sector). 
The similarity of results across rows shows that the sufficient-statistics approach accurately captures the welfare effects of interchange fees, even when accounting for consumer re-optimization and incomplete pass-through.

\end{table}

\clearpage

\begin{table}[H]
  \centering
  \caption{Effects of the Durbin Amendment}
  \label{tab:durbin}

  \begin{tabular}{lccccc}
\toprule
\multirow{2}{*}{Change in Consumption} & \multicolumn{5}{c}{Consumer Group} \\
& Cash & \makecell{Regulated \\ Debit} & \makecell{Exempt \\ Debit}  & \makecell{Basic \\ Credit}  & \makecell{Premium \\ Credit} \\
\midrule
Basis Points & 13 (4) & -34 (11) & 12 (3) & 12 (4) & 6 (3) \\
Dollars (\$Bn) & 1.9 (0.5) & -9.6 (3.0) & 3.3 (0.8) & 2.5 (0.8) & 1.8 (1.0) \\
\bottomrule
\end{tabular}
 
  \par\textit{Notes:} Table \ref{tab:durbin} shows estimated consumer welfare effects of the Durbin Amendment, which capped interchange fees on regulated debit cards. Welfare effects account for both capping baseline interchange fees and also eliminating fee heterogeneity. Standard errors for all entries are computed via a parametric bootstrap, described in Appendix \ref{subsec:appendix-bootstrap}.
\end{table}

\clearpage

\begin{table}[H]
  \centering
  \caption{The Rise of Premium Credit Cards}
  \label{tab:premiumization}

  \begin{tabular}{lccccc}
\toprule
\multirow{2}{*}{Change in Consumption} & \multicolumn{5}{c}{Consumer Group} \\
& Cash & \makecell{Regulated \\ Debit} & \makecell{Exempt \\ Debit}  & \makecell{Basic \\ Credit}  & \makecell{Premium \\ Credit} \\
\midrule
Basis Points of Consumption & -8 (1) & -9 (1) & -9 (1) & -8 (1) & 26 (2) \\
Dollars (billions) & -1.1 (0.1) & -2.6 (0.3) & -2.5 (0.2) & -1.7 (0.2) & 7.9 (0.7) \\
\bottomrule
\end{tabular}
 
  \par\textit{Notes:} Table \ref{tab:premiumization} summarizes the consumer-welfare effects of the increase in interchange fees on premium credit cards relative to a world in which premium cards charged the same interchange fees as basic cards. Welfare effects account for both the increase in overall fees and different levels of merchant heterogeneity in premium credit. Standard errors for all entries are computed via a parametric bootstrap, described in Appendix \ref{subsec:appendix-bootstrap}.
\end{table}

 
\appendix
\renewcommand{\thesubsection}{\Alph{section}.\arabic{subsection}}
\renewcommand{\thesubsubsection}{\Alph{section}.\arabic{subsection}.\arabic{subsubsection}}

\pagenumbering{arabic}
\renewcommand{\thepage}{A-\arabic{page}}
\renewcommand{\thetable}{A.\arabic{table}}
\renewcommand{\thefigure}{A.\arabic{figure}}
\setcounter{table}{0}  \setcounter{figure}{0}  

\clearpage{}


\section{Theory Appendix: Proofs and Details}\label{sec:appendix-theory}

This appendix formally demonstrates that the sufficient statistic in Equation~\ref{eq:sufficient-stat-formula} captures the first-order welfare effects of interchange fees, even when consumers are allowed to change their consumption in response to prices.

\subsection{Setup and Notation}\label{subsec:setup-notation}

Throughout the study, we index merchants (firms) by $j$, payment methods (card types) by $k$ and $l$, and sectors by $s$.
We define the following objects to characterize the equilibrium response of prices to interchange fee changes:

\begin{center}
\begin{tabular}{lll}
\toprule
\textbf{Symbol} & \textbf{Dimension} & \textbf{Economic Meaning} \\
\midrule
$\log p$ & $\mathbb{R}^J$ & Log retail prices \\
$q^k$ & $\mathbb{R}^J$ & Quantities purchased by type-$k$ consumers \\
$E^k$ & $\mathbb{R}^{J \times J}$ & Cross-price elasticities: $E^k_{ij} = \frac{\partial \log q_i^k}{\partial \log p_j}$ \\
$E$ & $\mathbb{R}^{JK \times J}$ & Stacked elasticity matrix $[E^1; \ldots; E^K]$ \\
$G$ & $\mathbb{R}^{J \times K}$ & Income semi-elasticities: $G^k_j = \frac{\partial \log q_j^k}{\partial y_k}$ \\
$F_j$ & $\mathbb{R}^{1 \times K}$ & Composition vector: $[(\tau_{j1} - \hat{\tau}_j)\mu_{j1}, \ldots, (\tau_{jK} - \hat{\tau}_j)\mu_{jK}]$ \\
$F$ & $\mathbb{R}^{J \times JK}$ & Block-diagonal composition matrix $\mathrm{blkdiag}(F_1, \ldots, F_J)$ \\
$\tau_{\max}$ & Scalar & Maximum interchange fee: $\max_{j,k} |\tau_{jk}|$ \\
$\omega_j$ & Scalar & Merchant $j$'s revenue share: $R_j / \sum_{j'} R_{j'}$ \\
\bottomrule
\end{tabular}
\end{center}

\noindent The composition matrix $F$ is the key object governing feedback effects. The term $(\tau_{jk} - \hat{\tau}_j)$ is the ``excess fee'' for payment method $k$ at merchant $j$---the deviation from that merchant's average fee $\hat{\tau}_j = \sum_k \mu_{jk} \tau_{jk}$. We use $\mathbb{E}_R[\cdot] = \sum_j \omega_j \cdot $ to denote expectations weighted by revenue shares, and normalize total revenue to 1 throughout.

\subsection{Assumption on Demand Elasticities}

We impose the following regularity condition on the elasticity matrix:

\begin{assumption}\label{assump:dominant-diagonal}
Elasticities are finite and satisfy a dominant diagonal condition.
There exist constants $M, m > 0$ such that for each $k$ and $j$:

\[
|E^k_{jj}| < M \quad \text{and} \quad \sum_{i \neq j} |E^k_{ji}| < m |E^k_{jj}|.
\]
\end{assumption}

This assumption ensures that own-price elasticities dominate cross-price elasticities and that all elasticities are bounded.
For example, in a CES system with $\sigma > 1$, we have $M=\sigma, m=1$.
Crucially, the bound $\|E\| \leq (1+m)M$ ensures the Neumann series in Lemma~\ref{lem:price-mapping} converges.

\begin{assumption}[Bounded Network Costs]\label{assump:positive-costs}
Card networks incur processing costs $c_k \in (0, \tau_{\max}]$ for each payment method $k$.\footnote{The bound $c_k \leq \tau_{\max}$ captures the requirement that networks are profitable on average---if processing costs exceeded interchange fees, networks would exit the market.}
This implies that consumer rewards satisfy $|f_k| \leq \tau_{\max}$, i.e., rewards are of order $O(\tau_{\max})$.

\end{assumption}

\subsection{Preliminary Results}\label{subsec:preliminary}

We establish two preliminary results used throughout the proofs.

\begin{lemma}[Spending Identity]\label{lem:spending-identity}
Under the normalization that total revenue equals 1, spending by type-$k$ consumers at merchant $j$ satisfies

\[
p_j q_{jk}^* = \mu_k^{-1} \omega_j \mu_{jk}.
\]
\end{lemma}

\begin{proof}
By definition, $\mu_{jk} = \mu_k q_{jk}^* / Q_j^*$ (the share of merchant $j$'s sales from payment method $k$) and $\omega_j = p_j Q_j^*$ (since total revenue is normalized to 1).
Therefore:

\[
\omega_j \mu_{jk} = p_j Q_j^* \cdot \frac{\mu_k q_{jk}^*}{Q_j^*} = \mu_k p_j q_{jk}^*.
\]
Dividing both sides by $\mu_k$ yields the result.

\end{proof}

\begin{lemma}[Composition Matrix Bound]\label{lem:F-bound}
The composition matrix satisfies $\|F\|_\infty \leq 2\tau_{\max}$.

\end{lemma}

\begin{proof}
Each entry of $F_j$ is $(\tau_{jk} - \hat{\tau}_j)\mu_{jk}$.
Since $\hat{\tau}_j = \sum_l \mu_{jl}\tau_{jl}$ is a convex combination of fees in $[-\tau_{\max}, \tau_{\max}]$, we have $|\tau_{jk} - \hat{\tau}_j| \leq 2\tau_{\max}$.
Since $\mu_{jk} \in [0,1]$ and $\sum_k \mu_{jk} = 1$, the row sum satisfies:

\[
\sum_k |(\tau_{jk} - \hat{\tau}_j)\mu_{jk}| \leq 2\tau_{\max} \sum_k \mu_{jk} = 2\tau_{\max}.
\]
\end{proof}

\subsection{Core Lemmas}\label{subsec:core-lemmas}

\paragraph{Reward Channel (First Moment).}
The following lemma characterizes the rewards that consumers receive from card networks.

\begin{lemma}[Rewards Approximation]\label{lem:rewards-approx}
In equilibrium, rewards for payment method $k$ satisfy

\[
f_k = \mu_k^{-1} \mathbb{E}_R[\mu_{jk}(\tau_{jk} - c_k)] + O(\tau_{\max}^2).
\]
\end{lemma}

\begin{proof}
By definition, $f_k = \frac{\sum_j p_j q_{jk}^*(\tau_{jk} - c_k)}{\mu_k(1 + f_k)}$, where we used that type-$k$ consumers spend $\mu_k(1+f_k)$.
Using $p_j q_{jk}^* = \mu_k^{-1} \omega_j \mu_{jk}$ gives $f_k(1 + f_k) = \mu_k^{-1} \mathbb{E}_R[\mu_{jk}(\tau_{jk} - c_k)]$.
Since $f_k = O(\tau_{\max})$, we have $f_k^2 = O(\tau_{\max}^2)$, yielding the result.

\end{proof}

\paragraph{Price Channel (Second Moment).}
The following lemma shows that composition feedback is second-order, so prices respond approximately one-for-one to direct fee changes.

\begin{lemma}[Price Pass-Through Mapping]\label{lem:price-mapping}
Suppose Assumption~\ref{assump:dominant-diagonal} holds and $2(1+m)M\tau_{\max} < 1$.
Then for any change in card-specific interchange fees $d\tau_{jk}$ (with direct effect on average fees $d\hat{\tau}_j = \sum_k \mu_{jk} d\tau_{jk}$) and $dy = O(\tau_{\max})$ change in incomes, the equilibrium change in log prices satisfies

\[
d \log p = d\hat{\tau} + O(\tau_{\max}),
\]
where $d\hat{\tau}$ is the vector of direct effects on average fees holding composition fixed.

\end{lemma}

\begin{proof}
The average fee $\hat{\tau}_j = \sum_k \mu_{jk} \tau_{jk}$ changes through the direct effect $d\hat{\tau}_j = \sum_k \mu_{jk} d\tau_{jk}$ and an indirect effect from composition changes $\sum_k \tau_{jk} d\mu_{jk}$.
Using the quotient rule for $d\mu_{jk}$, the indirect effect equals $\sum_k (\tau_{jk} - \hat{\tau}_j) \mu_{jk} d \log q_j^k$.
Combined with the demand system $d \log q = E \, d \log p + G \, dy$, this yields the fixed-point equation $d \log p = FE \, d \log p + FG \, dy + d\hat{\tau}$.

By Lemma~\ref{lem:F-bound}, $\|F\| \leq 2\tau_{\max}$, and by the dominant diagonal condition, $\|E\| \leq (1+m)M$.
Thus $\|FE\| < 1$ by hypothesis, so $(I - FE)^{-1}$ exists and the Neumann series converges.
All terms beyond $d\hat{\tau}$ are $O(\tau_{\max})$.

\end{proof}

\begin{remark}[Income Effects are Second-Order]\label{rem:income-effects}
In Lemma~\ref{lem:price-mapping}, the income change $dy_k = df_k$ is the change in rewards for type-$k$ consumers.
By Assumption~\ref{assump:positive-costs}, $|f_k| \leq \mu_k^{-1}\tau_{\max}$, so $dy = O(\tau_{\max})$.
Combined with $\|F\| = O(\tau_{\max})$ from Lemma~\ref{lem:F-bound}, the income feedback term $FG\,dy$ is $O(\tau_{\max}^2)$---second-order in fees.
This is why income effects do not appear in the first-order price mapping $d\log p = d\hat{\tau} + O(\tau_{\max})$.
\end{remark}

\begin{remark}[Composition Feedback is Second-Order]\label{rem:feedback}
The term $FE\,d\log p$ in the price equation represents composition feedback: Price changes cause consumers to reallocate across merchants, changing $\mu_{jk}$, which changes merchant costs.
Because $\|F\| = O(\tau_{\max})$ by Lemma~\ref{lem:F-bound}, this feedback contributes only $O(\tau_{\max}^2)$ to welfare (i.e., to the integrated welfare change from fees of size $O(\tau_{\max})$, consistent with the $O(\tau_{\max})$ error in the welfare derivative).
This justifies the sufficient statistic's focus on mechanical effects, as we do not need to solve for the full general equilibrium of prices and compositions.

\end{remark}

\subsection{Proof of Theorem 1}\label{subsec:proof-theorem1}

Below, we show that the first-order equilibrium welfare effect of a change in interchange fees for payment method $l$ on a type-$k$ consumer is given by aggregating the mechanical effects of changes in interchange fees on merchant-level prices.

\begin{proof}
By Roy's identity and the lemmas above:

\begin{align*}
\frac{1}{\lambda_k} \frac{dV_k}{dx_l}
&= \frac{df_k}{dx_l} - \sum_j q_{jk}^* \frac{dp_j}{dx_l} && \text{(Roy's identity)} \\
&= - \sum_j q_{jk}^* p_j \mu_{jl} \tfrac{\partial \tau_{jl}}{\partial x_l} + \mathbb{1}\{k=l\} \mu_k^{-1} \mathbb{E}_R\!\left[\mu_{jk} \tfrac{\partial \tau_{jl}}{\partial x_l}\right] + O(\tau_{\max}) && \text{(prices, rewards)} \\
&= \mu_k^{-1} \left( - \mathbb{E}_R\!\left[\mu_{jk} \mu_{jl} \tfrac{\partial \tau_{jl}}{\partial x_l}\right] + \mathbb{1}\{k=l\} \mathbb{E}_R\!\left[\mu_{jk} \tfrac{\partial \tau_{jl}}{\partial x_l}\right] \right) + O(\tau_{\max}) && \text{(spending identity)}
\end{align*}
where the second line uses Lemmas~\ref{lem:rewards-approx} and~\ref{lem:price-mapping}, and the third uses Lemma~\ref{lem:spending-identity}.

\end{proof}
\subsection{Proof of Fee Component Formulas}\label{subsec:fee-components}

We derive the fee component decomposition by applying Theorem 1 to specific counterfactual changes in the fee structure.
The welfare effect of any fee change equals the magnitude of the change multiplied by the derivative from Theorem 1.

\paragraph{Baseline Fees.}
Consider the effect of the baseline fee $\overline{\tau}_l$ relative to a zero-fee benchmark.
From the fee decomposition $\tau_{jl} = \overline{\tau}_l + \tau_{s_j,l} + b_j\delta_l$, we have $\frac{\partial \tau_{jl}}{\partial \overline{\tau}_l} = 1$ for all merchants.
Applying Theorem 1:

\begin{align*}
\frac{1}{\lambda_k}\frac{dV_k}{d\overline{\tau}_l}
&= \mu_k^{-1}\left( -\mathbb{E}_R[\mu_{jk}\mu_{jl}] + \mathbb{1}\{k=l\} \mathbb{E}_R[\mu_{jk}]\right) \\
&= \mu_k^{-1}\left( -\mathbb{E}_R[\mu_{jk}\mu_{jl}] + \mathbb{1}\{k=l\} \mu_k\right)
\end{align*}
where the second line uses $\mathbb{E}_R[\mu_{jk}] = \mu_k + O(\tau_{\max})$.
The total effect of the baseline fee is:

\[
A_{kl} = \overline{\tau}_l \times \frac{1}{\lambda_k}\frac{dV_k}{d\overline{\tau}_l} =\overline{\tau}_l \times \mu_k^{-1}\left(- \mathbb{E}_R[\mu_{jk}\mu_{jl}] + \mathbb{1}\{k=l\} \mu_k\right).
\]

\paragraph{Sector Adjustments.}
Consider the effect of sector adjustments $\tau_{sl}$ relative to uniform fees.
From the fee decomposition, $\frac{\partial \tau_{jl}}{\partial \tau_{sl}} = \mathbb{1}\{s_j = s\}$.
The key insight is that by the normalization $\sum_s \omega_{s} \mu_{sl} \tau_{sl} = 0$, where $\omega_{s} = \sum_{s_j = s} \omega_{j}$ is the sector's revenue weight and $\mu_{sl}$ is the revenue-weighted average share of spending on card $l$ among merchants in sector $s$. Thus, the reward term vanishes:
\[
\mathbb{E}_R[\mu_{jl} \tau_{s_j,l}] = \sum_s P_R(s_j=s) \mathbb{E}_R[\mu_{jl}|s_j=s] \tau_{sl} = 0.
\]
Therefore, when we apply Theorem 1 to changes in sector adjustments, the first-moment (reward) effect is zero, and only the second-moment (price) effect remains:

\[
\sum_s \tau_{sl} \times \frac{1}{\lambda_k}\frac{dV_k}{d\tau_{sl}} = \mu_k^{-1}\sum_s -\tau_{sl} P_R(s_j=s) \mathbb{E}_R[\mu_{jk}\mu_{jl}|s_j=s].
\]

\paragraph{Negotiated Rates.}
Consider the effect of the negotiated discount $\delta_l$ relative to standard rates.
From the fee decomposition, $\frac{\partial \tau_{jl}}{\partial \delta_l} = b_j$ (equals 1 for large merchants, 0 otherwise).
Applying Theorem 1:

\begin{align*}
\frac{1}{\lambda_k}\frac{dV_k}{d\delta_l}
= & \mu_k^{-1}\left( -\mathbb{E}_R[\mu_{jk}\mu_{jl} b_j] + \mathbb{1}\{k=l\} \mathbb{E}_R[\mu_{jk} b_j]\right) \\
= & \mu_k^{-1}P_R(b_j=1)\left( -\mathbb{E}_R[\mu_{jk}\mu_{jl}|b_j=1] + \mathbb{1}\{k=l\} \mathbb{E}_R[\mu_{jk}|b_j=1]\right).
\end{align*}

The total effect of the negotiated discount is therefore:
\[
B_{kl} = \delta_l \times \frac{1}{\lambda_k}\frac{dV_k}{d\delta_l} = \delta_l \times \mu_k^{-1}P_R(b_j=1)\left(-\mathbb{E}_R[\mu_{jk}\mu_{jl}|b_j=1] + \mathbb{1}\{k=l\} \mathbb{E}_R[\mu_{jk}|b_j=1]\right).
\]

\paragraph{Welfare Under Imperfect Pass-Through.}
The sensitivity analysis in Table \ref{tab:sensitivity-analysis} relaxes the assumption of full pass-through by introducing pass-through rates $\rho^P \in (0,1]$ for prices and $\rho^R \in (0,1]$ for rewards.
Under imperfect pass-through, the welfare formula becomes
\[
\frac{1}{\lambda_k}\frac{dV_k}{dx_l}
\approx \mu_k^{-1}\left(
-\rho^P\,\mathbb{E}_R\!\left[\mu_{jk}\mu_{jl}\tfrac{\partial\tau_{jl}}{\partial x_l}\right]
+ \rho^R\,\mathbb{1}\{k=l\}\,\mathbb{E}_R\!\left[\mu_{jk}\tfrac{\partial\tau_{jk}}{\partial x_k}\right]
\right).
\]
When both pass-through rates equal 0.7, redistribution dampens proportionally.
When only price pass-through falls to 0.7, while rewards pass-through remains at 1, credit card users benefit more than in the baseline because they receive full rewards while bearing only partial price increases.
Conversely, when only rewards pass-through falls to 0.7, credit card users benefit less because they face full price increases while receiving only partial rewards.

 
\clearpage{}


\section{Data Cleaning and Construction of Core Datasets}

This section describes the cleaning procedures applied to five key datasets that form the backbone of our empirical analysis: (1) the aggregate time series, (2) the Clover merchant data, (3) the merchant cross-section from the settlement dataset, (4) the Diary of Consumer Payment Choice, and (5) the MRI survey data.

\subsection{Aggregate Payment Time Series}

We aggregate Fiserv settlement data from 2006--2022 to provide aggregate measures of payment instrument usage over time.
We construct this data by aggregating the raw merchant-level information over time into a transaction code by month panel.
When showing time series, we drop transaction codes associated with negotiated fees (e.g., large merchants).
By dropping large firms, we remove the influence of entry and exit of large merchants over time, which can create mechanical changes in average fees and volumes.
We also drop transaction codes with less than zero net interchange in a given month.
Interchange codes can be negative for a month due to refunds or adjustments.

\subsection{Clover Dataset}

The Clover data provide establishment-level records from a sample of merchants spanning 2019--2022 that includes both card and cash transactions.
We first aggregate establishment-level transactions into annual cash, check, debit, and credit card volumes.
We then drop any merchants with negative sales, as those reflect the influence of net refunds and chargebacks.
The Clover data provide our only measure of cash usage at the merchant level.

\subsection{Settlement Dataset}

The settlement cross-section provides establishment-level aggregates of card transactions and interchange fees for 2022.
We focus on 2022 to ensure comparability with the Clover data.
We construct a cross-section of merchants that can be consistently classified by sector and linked to firm identifiers.
Cleaning proceeds in several steps:

\begin{enumerate}
    \item Merchants with missing or inconsistent MCC codes are dropped.
    \item Financial transactions (such as payment app transfers and ATM withdrawals) are removed to focus on true retail purchases.
    \item We then collapse the payment volumes and interchange fees to the establishment level. Note that multiple establishments may belong to the same firm.
    \item For the sufficient statistics analysis, we further collapse the payment volumes and fees to the firm level by aggregating across all establishments belonging to the same firm.
\end{enumerate}

We group MCC codes into six broad sectors: (i) Grocery, (ii) Dining, (iii) Gas Stations, (iv) General Retail, (v) Travel, and (vi) Other.
We choose these categories because they are the largest broad categories of merchants that are also present in the Diary of Consumer Payment Choice.
Grocery includes warehouse stores, discount stores, and supermarkets.
Dining includes full-service and quick-service restaurants, as well as bars.
General Retail includes department stores, specialty retailers, discount stores, pharmacies, and convenience stores.
Travel includes airlines, hotels, car rentals, and travel agencies.
Other includes all remaining MCC codes, such as healthcare, utilities, and entertainment.

Because the settlement dataset contains only card transactions, we impute cash usage using a statistical model estimated on the Clover data.
The modeling approach, validation of the extrapolation procedure, and detailed results are described in Appendix~\ref{sec:appendix-cash-model}.

\subsection{Consumer Surveys: Diary of Consumer Payment Choice (DCPC) and MRI}

We use two consumer survey datasets to analyze household payment preferences and card ownership patterns.
This allows us to move beyond redistribution across consumers of different payment methods and instead study how interchange fees redistribute across households with different income levels.

The Diary of Consumer Payment Choice collects individual information on demographics, payment preferences, and transaction diaries.
We pool the 2022--2023 panels to maximize the sample size for our analysis.
The data on demographics and payment preferences allow us to model the correlation between household income and preferred payment method.
We define preferred payment methods as consumers' preferred method for in-person transactions.
We focus on those who report cash, debit, or credit, which essentially captures all consumers.
We also use the DCPC transaction diaries to measure the distribution of payment methods across different sectors.
This helps us validate the sector-level payment patterns observed in the settlement data in Figure \ref{fig:data-representative-sector}.

The MRI survey is a large, nationally representative consumer survey that links demographics and brand preferences.
Among the many variables it collects are household demographics, card ownership, and bank use.
This allows us to separate households into using exempt versus regulated debit, and also premium versus regular credit cards.
We use the 2022 wave which covered around 50,000 households.

\begin{itemize}
    \item We code a household as owning a credit (or debit) card if it reports owning at least one of that type. We also track multiple ownership and identify network-level ownership (Visa, Mastercard, Amex, Discover). We exclude households reporting neither credit nor debit spending and not preferring cash (17\% of the MRI households), as they provide no information on payment preferences.
    \item To identify regulated versus exempt debit, we group financial institutions into two categories. \emph{Small banks} include credit unions, community banks, and Chime. \emph{Big banks} include large national and regional banks (Chase, Citi, U.S.Bank, Wells Fargo, Bank of America, Ally, BB\&T, BBVA, Capital One, Citizens, Fifth Third, HSBC, Key, PNC, Regions, SunTrust, TD). We split debit into exempt and regulated in proportion to the number of institutions of each type reported. For example, if a household reports using two small and one big bank, 2/3 of debit spend is assigned to exempt and 1/3 to regulated.
    \item We define premium credit cards based on the types of cards people hold. We classify Visa/Mastercard Standard, Gold, and Platinum as \emph{regular} cards. We classify Visa Signature, Mastercard World Elite, and American Express as \emph{premium} cards. If a household owns any premium card, all credit spending is assigned to the premium category; otherwise, spending is assigned entirely to regular cards.
\end{itemize}

With these two datasets, for discrete income bin categories, we can compute the distribution of payment preferences across debit, credit, and cash from the DCPC, and then we split debit into regulated versus exempt and credit into premium versus regular using the MRI data.
Table~\ref{tab:dcpc-summary-stats} shows the summary statistics for both datasets.

\begin{table}[H]

  \centering
  \caption{Summary Statistics: Diary of Consumer Payment Choice (DCPC) and MRI Survey}
  \label{tab:dcpc-summary-stats}
  \begin{tabular}{lrrr}
\toprule
 & N & Mean & SD \\
\midrule
\multicolumn{4}{l}{\textit{Panel A: DCPC}} \\
\quad Household income (\$) & 8,547 & 88,505 & 74,497 \\
\quad Prefers debit (\%) & 8,547 & 42 &  \\
\quad Prefers credit (\%) & 8,547 & 38 &  \\
\quad Prefers cash (\%) & 8,547 & 20 &  \\

\multicolumn{4}{l}{\textit{Panel B: MRI}} \\
\quad Household income (\$) & 51,697 & 93,369 & 76,509 \\
\quad Prefers debit (\%) & 51,697 & 31 &  \\
\quad Prefers credit (\%) & 51,697 & 42 &  \\
\quad Prefers cash (\%) & 51,697 & 29 &  \\
\quad Exempt debit share (\%) & 15,686 & 36 &  \\
\quad Premium credit share (\%) & 22,016 & 45 &  \\
\bottomrule
\end{tabular}
 
  \par\textit{Notes:} The table shows summary statistics for respondents in the DCPC and MRI survey. Panel A shows demographics and payment preferences from the DCPC. Payment preferences are defined by consumers' responses to their preferred in-person payment method. Panel B shows card ownership and bank use from the MRI survey. Cash preference is defined by whether consumers report that they use cash whenever possible. Conditional on not reporting a cash preference, consumers are assigned to credit or debit based on whether they spend more on each card type. Regulated versus exempt debit is defined based on the types of financial institutions used by the household (see text for details). Premium versus regular credit is defined based on the types of credit cards owned by the household (see text for details). Statistics are reported with sampling weights.
\end{table}

Relative to the aggregate payment patterns in the settlement data (approximately 11\% cash, 45\% debit, 44\% credit), consumers in the DCPC report a higher preference for cash (around 20\% in the survey), and consumers in the MRI survey report a higher preference for debit cards from large banks and for regular credit cards.
To address this bias, we proportionally re-weight shares of payment methods within each income bin to match the aggregate shares observed in the settlement data.
This re-weighting preserves the correlation between income and payment preferences observed in the surveys while ensuring that the overall payment method shares align with the settlement data.
This proportional rescaling is equivalent to a logit adjustment: we solve for method-specific constants that shift aggregate predicted shares to match the settlement data while preserving relative preferences across income bins.





\subsection{Validation Against External Industry Benchmarks}

To assess the external validity of our datasets, we compare key metrics from the Fiserv data against independent measures from the Nilson Report, a leading industry trade publication that compiles payment statistics across the U.S. payment system.
Figure~\ref{fig:fiserv-nilson-representativeness} compares transaction sizes, credit card shares, average fees, and paper payment shares.
The Fiserv data closely track the Nilson benchmarks across all four dimensions.
Credit card transaction sizes in Fiserv are approximately 20\% lower than the Nilson Report, reflecting our exclusion of business credit cards; debit card transaction sizes are only about 5\% lower.
The credit-debit mix, fee levels, and cash shares all align closely with external estimates, validating both our settlement data and our Clover-based extrapolation procedure for cash usage.

\clearpage

\begin{figure}[H]
  \centering
  \caption{Representativeness of Fiserv Data: Validation Against Nilson Report Industry Benchmarks}
  \label{fig:fiserv-nilson-representativeness}

\begin{subfigure}[t]{0.48\textwidth}
    \centering
    \includegraphics[width=\linewidth]{tmp_images/___fiserv_output_graphs_fiserv_representativeness_tx_size_pdf.png}
    \caption{Average Transaction Size}
    \label{fig:fiserv-nilson-tx-size}
  \end{subfigure}\hfill
  \begin{subfigure}[t]{0.48\textwidth}
    \centering
    \includegraphics[width=\linewidth]{tmp_images/___fiserv_output_graphs_fiserv_representativeness_credit_share_pdf.png}
    \caption{Credit Card Share of Card Payments}
    \label{fig:fiserv-nilson-credit-share}
  \end{subfigure}

  \vspace{0.6em}
  \begin{subfigure}[t]{0.48\textwidth}
    \centering
    \includegraphics[width=\linewidth]{tmp_images/___fiserv_output_graphs_fiserv_representativeness_avg_fee_pdf.png}
    \caption{Average Fee per Transaction}
    \label{fig:fiserv-nilson-avg-fee}
  \end{subfigure}\hfill
  \begin{subfigure}[t]{0.48\textwidth}
    \centering
    \includegraphics[width=\linewidth]{tmp_images/___fiserv_output_graphs_fiserv_representativeness_paper_share_pdf.png}
    \caption{Share of Paper-Based Payments}
    \label{fig:fiserv-nilson-paper-share}
  \end{subfigure}

  \par\textit{Notes:} Validation of Fiserv data representativeness against Nilson Report industry benchmarks. Panel (a): Average transaction size in the Fiserv data is measured as total dollar volume divided by number of transactions.
Panel (b) compares the dollar share of V/MC volume on credit cards.
Panel (c): Fees in the Nilson report represent interchange fees as well as additional acquirer fees.
Interchange in the Fiserv data excludes acquirer fees.
Panel (d) plots the share of all transactions via cash/checks in the Clover data against cash as a share of cash + card transactions from the Nilson Report.
We exclude checks in the Nilson report because they are mainly for non-retail sectors that are not well represented in the Fiserv data.

\end{figure}

\clearpage

\clearpage{}


\section{Cash Extrapolation Model: From Clover to Settlement}
\label{sec:appendix-cash-model}

The Clover dataset sheds light on the distribution of cash usage across merchants, which is not directly observable in the settlement dataset.
However, Clover merchants are not representative of the composition of the broader economy.
We therefore build a model to extrapolate cash usage from Clover to the broader settlement sample.
This section details the modeling approach, estimation results, and validation exercises.

\subsection{Overview}

We use a two-part model to extrapolate cash usage from Clover to the broader settlement sample.

\begin{enumerate}
    \item A logistic regression for the probability that the share of cash transactions at a merchant exceeds 2\% (extensive margin).
    \item A linear regression for the log of the conditional cash share among merchants with greater than 2\% share (intensive margin).
\end{enumerate}
We use a two-stage model because the distribution of cash use in Clover data has a large mass point at zero and is highly skewed among positive values.

An important motivation for using the regression model is that Clover merchants are very different from the merchants in the broader Fiserv settlement sample.
They are much smaller and much more likely to be in the restaurant and retail sectors.
Thus, we cannot simply take the average cash share in Clover and apply it to the settlement data.
Instead, we need to model how cash usage varies with observable merchant characteristics and then use that model to predict cash usage in the settlement data.

\subsection{Methodology and Results of Cash Prediction Model}

We find that our two-stage approach effectively captures the variation in cash usage across merchants in the Clover data.
The resulting model matches key features of the distribution of cash use in the Clover data and is consistent with several pieces of external evidence.

\subsubsection{Model Specification and Estimation Results}

We estimate both parts of the cash extrapolation model on the Clover sample.
The key predictors in both stages include the average ticket size of card transactions, merchant-level share of online transactions, firm size, sector fixed effects, and local demographic controls.
The strongest predictors of cash usage are average ticket size, e-commerce participation, firm size, and the share of card transactions on debit cards.
The model also includes demographic controls (population density, median household income, education, age, race/ethnicity) and industry fixed effects.

\paragraph{Ticket Size.} Higher ticket sizes are strongly associated with lower cash usage on both margins. Figures~\ref{fig:cash_prob_ticket} and \ref{fig:cash_condl_ticket} illustrate this univariate relationship in the Clover data: cash usage declines sharply as average ticket size increases.

\paragraph{E-commerce Status.} E-commerce merchants have substantially lower cash usage, as expected. Figures~\ref{fig:cash_prob_ecomm} and \ref{fig:cash_condl_ecomm} show the difference in cash usage on the extensive and intensive margins, respectively, between brick-and-mortar and e-commerce merchants.

\paragraph{Firm Size.} Larger firms tend to have higher probability of accepting cash but lower conditional cash shares. Figures~\ref{fig:cash_prob_size} and \ref{fig:cash_condl_size} display how the extensive and intensive margins of cash usage vary with firm size.

\paragraph{Debit Share.} For most of the sample, higher debit share is associated with higher cash usage on both margins. This is intuitive given that debit cards are a closer substitute for cash than they are for credit cards. This pattern reverses at very high debit shares (above roughly 80\%), where cash usage declines with debit share. This non-linear relationship is captured by including both debit share and its square in the regression. Figures~\ref{fig:cash_prob_debit} and \ref{fig:cash_condl_debit} illustrate this relationship.

\begin{figure}[ht]
    \centering
    \caption{Probability of Positive Cash by Merchant Characteristics (Extensive Margin)}
    \label{fig:cash_extensive}

\begin{subfigure}[t]{0.48\textwidth}
        \centering
        \includegraphics[width=\linewidth]{tmp_images/___fiserv_output_graphs_clover_cash_prob_pos_ticket_pdf.png}
        \caption{By Average Ticket Size}
        \label{fig:cash_prob_ticket}
    \end{subfigure}\hfill
    \begin{subfigure}[t]{0.48\textwidth}
        \centering
        \includegraphics[width=\linewidth]{tmp_images/___fiserv_output_graphs_clover_cash_prob_pos_ecomm_pdf.png}
        \caption{By E-commerce Status}
        \label{fig:cash_prob_ecomm}
    \end{subfigure}

    \vspace{0.6em}

\begin{subfigure}[t]{0.48\textwidth}
        \centering
        \includegraphics[width=\linewidth]{tmp_images/___fiserv_output_graphs_clover_cash_prob_pos_size_pdf.png}
        \caption{By Firm Size}
        \label{fig:cash_prob_size}
    \end{subfigure}\hfill
    \begin{subfigure}[t]{0.48\textwidth}
        \centering
        \includegraphics[width=\linewidth]{tmp_images/___fiserv_output_graphs_clover_cash_prob_pos_pdf.png}
        \caption{By Debit Share}
        \label{fig:cash_prob_debit}
    \end{subfigure}

    \par\textit{Notes:} Figure shows the probability that a merchant has more than 2\% of transactions in cash, by merchant characteristics in the Clover data. Panel (a): Higher ticket sizes are associated with lower probability of cash usage. Panel (b): E-commerce merchants are less likely to have positive cash usage. Panel (c): Larger firms are more likely to accept cash. Panel (d): Higher debit share is generally associated with higher probability of cash usage.
\end{figure}

\begin{figure}[ht]
    \centering
    \caption{Conditional Cash Share by Merchant Characteristics (Intensive Margin)}
    \label{fig:cash_intensive}

\begin{subfigure}[t]{0.48\textwidth}
        \centering
        \includegraphics[width=\linewidth]{tmp_images/___fiserv_output_graphs_clover_cash_condl_share_ticket_pdf.png}
        \caption{By Average Ticket Size}
        \label{fig:cash_condl_ticket}
    \end{subfigure}\hfill
    \begin{subfigure}[t]{0.48\textwidth}
        \centering
        \includegraphics[width=\linewidth]{tmp_images/___fiserv_output_graphs_clover_cash_condl_share_ecomm_pdf.png}
        \caption{By E-commerce Status}
        \label{fig:cash_condl_ecomm}
    \end{subfigure}

    \vspace{0.6em}

\begin{subfigure}[t]{0.48\textwidth}
        \centering
        \includegraphics[width=\linewidth]{tmp_images/___fiserv_output_graphs_clover_cash_condl_share_size_pdf.png}
        \caption{By Firm Size}
        \label{fig:cash_condl_size}
    \end{subfigure}\hfill
    \begin{subfigure}[t]{0.48\textwidth}
        \centering
        \includegraphics[width=\linewidth]{tmp_images/___fiserv_output_graphs_clover_cash_condl_share_pdf.png}
        \caption{By Debit Share}
        \label{fig:cash_condl_debit}
    \end{subfigure}

    \par\textit{Notes:} Figure shows the conditional cash share among merchants with more than 2\% cash usage in the Clover data, by merchant characteristics. Panel (a): Higher ticket sizes are associated with lower cash usage. Panel (b): E-commerce merchants have substantially lower cash usage than brick-and-mortar merchants. Panel (c): Larger firms tend to have lower conditional cash shares. Panel (d): Higher debit share is generally associated with higher cash usage.
\end{figure}

\subsubsection{Model Performance}

The model does an adequate job of predicting both the extensive and intensive margin of cash use.
The extensive margin model achieves an AUC of 0.82, indicating good discriminative ability in predicting which merchants have more than 2\% of their transactions on cash.

The intensive margin model achieves a correlation of 0.52 between predicted and actual log conditional cash shares.
Figure~\ref{fig:condl_distribution_fit} shows the fit of the predicted conditional cash share distribution to the actual distribution in the Clover data, validating that the model captures the relevant variation in cash usage.

\begin{figure}[ht]
    \centering
    \caption{Conditional Cash Share: Model Fit vs.\ Actual Distribution}
    \label{fig:condl_distribution_fit}

    \includegraphics[width=0.7\textwidth]{tmp_images/___fiserv_output_graphs_clover_condl_distribution_fit_pdf.png}

    \par\textit{Notes:} Figure compares the predicted conditional cash share distribution from the model to the actual distribution in the Clover data. The close fit validates that the model captures the relevant variation in cash usage across merchants.
\end{figure}

\subsubsection{Aggregate Validation against Nilson Reports and Homescan Data}

Our sufficient-statistics approach relies on accurately estimating the second moments of payment shares across merchants.
Therefore, it is important that our cash extrapolation model accurately captures the overall level of cash usage in the economy and its variation across merchants.
After we fit the cash prediction model on the Clover data, we apply it to the settlement data to predict cash usage across the settlement data. When extrapolating to the settlement data, we cap the log card sales variable at the 99th percentile of the Clover distribution to avoid extreme extrapolations along the firm size dimension.

We find that our model accurately recovers both the overall level of cash usage and its variation across merchants in the settlement data.
According to the 2022 Nilson Report, the average use of cash in retail transactions is around 10\% \citep{Nilson2023a}.
Our model predicts an average cash share of 10\% across the settlement data, closely matching the Nilson benchmark.
Appendix E of \citet{Wang2025} reports the distribution of cash share across Homescan merchants, and finds a standard deviation of approximately 8-10 percentage points.
Our model predicts a standard deviation of 9 percentage points across the settlement data, again closely matching the Homescan benchmark.





































\clearpage{}


\section{Technical Details for Structural Model and Parametric Bootstrap}\label{sec:appendix-calibrated}

This appendix provides technical details for the structural model presented in Section \ref{sec:redistribution}.
The main text presents the model specification, calibration strategy, and welfare comparison.
This appendix contains the calibration goodness-of-fit diagnostics and detailed parameter values.

\subsection{Detailed Parameterization}

Our parameterization treats each market $m$ as featuring $j^{*}$ large firms and a fringe of monopolistically competitive firms that vary in the composition of their payments.
We parameterize the taste shifters as
\[
\log a_{jksm}=\log\frac{\overline{a}_{s}}{j^{*}}\times I\left\{ j\leq j^{*}\right\} +\log\frac{1-\overline{a}_{s}}{J-j^{*}}\times I\left\{ j>j^{*}\right\} +\log z_{jksm}
\]
The special intercept $\overline{a}_{s}$ for the first $j^{*}$ firms allows us to model these firms as large.
The $\log(J-j^{*})$ term is a normalization so that the sum of the small firms has a revenue share of approximately $1-\overline{a}_{s}$.

We model variation in payment composition by allowing consumers with different payment preferences to vary in their preferences over merchants.
We then set the payment-type taste adjustments to $\left(z_{j1sm},\dots,z_{jKsm}\right)\sim N\left(\overline{z}_{js},\Sigma_{js}\right)$.
For small firms with $j>j^{*}$, we set $\overline{z}_{js}=\mathbf{0}$ and impose a common covariance $\Sigma_{s}$.
For large firms with $j\leq j^{*}$, we impose a common mean and variance, $\overline{z}_{js}=\overline{z}_{s}^{*}$, $\Sigma_{js}=\Sigma_{s}^{*}$.

\subsection{Welfare Formulas}

To evaluate the welfare effects of interchange fees, we compute average compensating variation for each consumer type across the $M$ markets.
Since networks mechanically pass through fees to rewards and do not earn economic profits in our model, the welfare analysis focuses on consumers and merchants.
The key welfare metric for consumers is the change in the sectoral price indices they face, which captures both the direct effect of higher retail prices (from merchant fees) and the indirect effect of rewards.

Formally, in this framework, network profits are constant.
Average compensating variation for consumers of type $k$, expressed as a percentage of baseline expenditure, is
\[
\Delta\log W_{k}=\frac{1}{M}\sum_{m}\sum_{s}\alpha_{ks}\Delta\log P_{ksm}
\]

\subsection{Calibration Details}

In practice, given a set of parameters, we simulate $M=200$ markets, with each sector in each market containing $100$ firms, and then compute the resulting equilibrium.
We match the following moments:

\begin{itemize}
\item We set $\alpha_{ks}$ to match the share of spending of each payment
method in each sector.
This mechanically ensures that we match the average composition of payments in each sector.

\item In the merchandise and grocery sectors, we set $\overline{a}$ to match the revenue share of firms with more than \$1 billion in revenue. In the other sectors, we set $\overline{a}=0.49$ to match the merchandise sector.
\item We set $j^*$ to equal the same value across all sectors. In the complete pass-through calibration we set $j^*=50$, whereas for incomplete pass-through, we set $j^*=2$. Intuitively, $j^*$ governs the number of firms with strategic complementarities in each market. It is important to match the share of revenue at firms with more than \$1 billion in sales in grocery and merchandise because that determines the fee discounts given by the networks.
\item We calibrate $\Sigma_{s}$ to match the revenue-weighted second-moment matrices.
\item We set $\overline{z}_{s}^{*}$ and $\Sigma_{s}^{*}$ to match the revenue-weighted average composition of payments of firms with more than \$1 billion in revenue in those sectors.
\item We set $\sigma=5$. This is between the median across-firm and within-firm elasticities reported by \citet{Hottman.etal2016}. In any case, the aggregate elasticity does not affect the validity of the sufficient statistic approach.
\item We obtain the components of the fee $\tau_{jk}$ from regressions of firm-level fees on payment method $k$ as a share of sales on payment method $k$ on sector indicators and an indicator for being a large firm in the merchandise or grocery sectors.
\end{itemize}

\subsection{Goodness-of-Fit Diagnostics}

Figure \ref{fig:calibration-fit} demonstrates the quality of the calibration fit.
Panel \ref{fig:calibration-moments} shows that the model successfully replicates the revenue-weighted second moments of payment shares across sectors---the key statistics that govern the extent of cross-subsidization in our sufficient-statistics framework.
The close alignment between model-generated moments (y-axis) and data moments (x-axis) confirms that the calibration captures the observed covariance structure of payment methods across merchants.
Panel \ref{fig:markups-by-sector} shows that the model generates realistic markups by sector, which is important for ensuring that merchant pricing behavior is empirically grounded.
Finally, Panel \ref{fig:calibration-passthrough} regresses the change in each firm's average fees as a result of capping fees to zero in the counterfactual against the change in the firm's average price with market fixed effects. The slopes represent pass-through rates of idiosyncratic costs implied by the model. Whereas small firms in the structural calibration fully pass through own costs, large firms only pass through around half.
Together, these panels confirm that the calibrated model provides a credible quantitative laboratory for evaluating the welfare effects of interchange fee policies.

\begin{figure}[H]
  \centering
  \caption{Goodness of Fit: Model Calibration Results}
  \label{fig:calibration-fit}

  \begin{subfigure}[t]{0.9\textwidth}
    \centering
    \includegraphics[width=\linewidth]{tmp_images/___output_graphs_calibration_fit_by_sector_pdf.png}
    \caption{Second Moments by Sector}
    \label{fig:calibration-moments}
  \end{subfigure}

  \begin{subfigure}[t]{0.4\textwidth}
    \centering
    \includegraphics[width=\linewidth]{tmp_images/___output_graphs_markups_by_sector_pdf.png}
    \caption{Markups by Sector}
    \label{fig:markups-by-sector}
  \end{subfigure}
  \begin{subfigure}[t]{0.4\textwidth}
    \centering
    \includegraphics[width=\linewidth]{tmp_images/___output_graphs_firm_fe_passthrough_by_sector_extreme_pdf.png}
    \caption{Pass-through Regressions by Sector}
    \label{fig:calibration-passthrough}
  \end{subfigure}

  \par\textit{Notes:} Figure \ref{fig:calibration-fit} demonstrates that the model calibration successfully matches the targeted moments from the data. Panel \ref{fig:calibration-moments} shows the matching of revenue-weighted second moments for both large and small firms across sectors, a key component of the model's calibration strategy. Panel \ref{fig:markups-by-sector} displays markups by sector, showing the model's ability to match empirical markup patterns. Panel \ref{fig:calibration-passthrough} illustrates pass-through regressions by sector, providing further validation of the model's fit. Together, these figures provide visual evidence that the calibration is working as intended.

\end{figure}

\subsection{Parametric Bootstrap Procedure}
\label{subsec:appendix-bootstrap}

We compute standard errors using a parametric bootstrap that combines two sources of uncertainty: uncertainty in revenue-weighted payment moments, and uncertainty in fee parameters.

For the revenue-weighted moments, we must model both the distribution of payment shares across firms and the distribution of firm size.

To match the firm size distribution by sector, within each sector, we draw firm revenues from log-normal distributions calibrated to the empirical size distribution.
For the grocery and merchandise sectors, we take care to draw sufficiently many firms with more than \$1 billion in revenue by modeling the number of such firms as a Poisson random variable with mean equal to the observed number of such firms in that sector, and then modeling firm size above \$1 billion using a truncated log normal distribution.

To match the distribution of payment shares within each sector, we draw firm-level payment shares from a multivariate distribution calibrated to match the observed mean shares and their full covariance structure.
We do this by following the log-normal specification used in the structural calibration, which allows us to model the covariance matrix of shares while ensuring that the shares lie on the simplex.
For the grocery and merchandise sectors, we again take care to separately model the distribution of payment shares for firms with more than \$1 billion in revenue, drawing these from a distinct multivariate distribution calibrated to match the observed mean shares and covariance matrix for such firms.

For the fee parameters, we account for the fact that although our transaction-level regressions contain millions of observations, the effective sample size is much smaller due to the skewed firm size distribution.
We then draw fee coefficients from their asymptotic multivariate normal distribution.

Each bootstrap iteration combines a new draw of revenue-weighted moments with a new draw of fee parameters to compute the redistribution effects.
Standard errors are the standard deviation across iterations.

\clearpage

\clearpage{}


\section{Causal Evidence on Interchange Fee Pass-Through\label{sec:appendix_durbin}}

This appendix provides empirical evidence that small restaurants pass interchange fees through to higher retail prices.
We use the 2011 Durbin Amendment as a natural experiment to estimate the causal effects of interchange fees on merchant prices and sales among restaurants.
We focus on restaurants because the Durbin Amendment created econometrically convenient cross-sectional variation in this sector.\footnote{ Other sectors experienced more uniform fee changes under the Durbin Amendment.
According to Visa's public interchange fee schedules \citep{Visa2010}, the change in fixed fees was much smaller for retailers and gas stations (from 15 cents to 22 cents), whereas there was no change in the linear fee for grocers.
The restaurant sector thus represents one of the few settings with substantial heterogeneity in treatment intensity. } These findings help support our assumption that firms pass on interchange fees into retail prices and complement the findings from \citet{Amiti.etal2019} that sector-level shocks like interchange fees are passed into retail prices.

\subsection{The Heterogeneous Effects of the Durbin Amendment}

One challenge in studying the real effects of interchange fees is that it requires both granular data---at the merchant level---on interchange fees and merchant outcomes and exogenous variation in interchange fees.
Using our detailed merchant-level data, we exploit heterogeneity in the impact of the Durbin Amendment across merchants to show that interchange fees affect the allocation of consumption and retail prices.

The 2011 Durbin Amendment, passed as part of the Dodd-Frank Act, imposed a binding cap on debit card interchange fees for banks with over \$10 billion in assets, a policy that had highly heterogeneous effects across merchants.
As illustrated in Figure \ref{fig:interchange_ts}, the Durbin Amendment had a substantial impact on debit interchange fees while having no effect on credit interchange.

The amendment created sharp, plausibly exogenous variation in interchange fee savings across merchants, which was driven by three forces.
First, only merchants with substantial sales via debit cards saw cost savings.
Second, because the cap includes a fixed 22-cent component, the percentage reduction in fees was larger for merchants with larger average transaction sizes.\footnote{ An inspection of Visa's public interchange fee schedules before and after implementation of the policy suggest that most regulated debit fees bunched at the cap.
In principle, the Durbin Amendment could have forced the networks to cut debit fees for the largest merchants even more to incentivize these merchants to stay on the Visa and Mastercard networks.
However, Figure \ref{fig:price-disc-size} shows regulated debit fees are similar at both large and small firms, suggesting this additional negotiation channel was not a major force. } Third, the amendment's impact varied by region: Only merchants in areas with large, regulated banks experienced fee changes, while merchants in areas without such banks were unaffected, providing a natural control group.
We use our detailed merchant-level data to measure these forces precisely.

We then construct a panel of annual restaurant card sales from October 2009 to October 2012.
Table \ref{tab:durbin-summary} characterizes this sample, showing substantial variation in transaction sizes and regulated bank exposure---the two key sources of identification.

We can then use this panel to illustrate the heterogeneous effects of the Durbin Amendment across restaurants.
Figure \ref{fig:durbin-het-ticket} binscatters average firm-level debit interchange rates, $\tau_j$, against the average pre-regulation transaction size of debit cards, $t_j$.
We measure firm-level interchange fees as the dollar value of interchange fees divided by the total dollar volume of transactions, and we measure transaction size as the total dollar volume of debit transactions divided by the total number of debit transactions.
By comparing the two curves, we see that the regulation increased interchange fees for restaurants with low ticket sizes and decreased them for restaurants with large ticket sizes.
This change occurred because the Durbin Amendment modified the structure of debit interchange for restaurants from a large linear fee with a low fixed fee ($1.19\%$ of transaction value + $10$ cents) to a low linear fee with a high fixed fee ($0.05\%$ of transaction value + $22$ cents) \citep{Visa2010}.

Figure \ref{fig:durbin-het-regulated} binscatters average debit interchange fees against the post-Durbin share of debit transactions on regulated cards at the ZIP code level, $r_{z(j)}$, both before and after the regulation.
Using post-Durbin data, we calculate the share of debit card transactions at restaurants in each ZIP code conducted on regulated cards.\footnote{ We measure the regulated bank share at the ZIP-code level rather than merchant level to address endogeneity concerns: merchant-level bank shares measured after the regulation may reflect endogenous changes in merchant customer composition, while ZIP-code level variation provides plausibly exogenous geographic variation in treatment intensity. } After the Durbin Amendment, average interchange rates declined sharply in areas with a higher share of regulated banks.

\subsection{Estimation}

We implement an instrumental variables design that compares restaurants that vary in their expected gains from the Durbin Amendment based on their pre-regulation characteristics across areas with differential exposure to regulated banks.
We estimate the following:

\begin{equation}
\Delta y_{j}=\beta\Delta\tau_{j} + \gamma X_j + \delta_{z(j)} + \epsilon_{j}
\end{equation}
where $\Delta y_{j}$ is the change in log sales at merchant $j$ in the year following the Durbin Amendment (October 2011 to October 2012), $\Delta\tau_{j}$ measures the average change in credit and debit interchange fees (expressed as percentage points) at merchant $j$ following the Durbin Amendment, $X_j$ are firm-level controls, and $\delta_{z(j)}$ are ZIP-code fixed effects.
The coefficient $\beta$ measures the percentage change in sales for each percentage-point change in fees.\footnote{ We model pass-through using a standard log-linear specification common in the tax literature.
The standard log-log pass-through regression regresses $\log P\left(1+\tau\right)$ on $\log\left(1+\tau\right)$ \citep{Butters.etal2022}.
Since interchange fees are small relative to prices, this is nearly identical to our specification on the right-hand side, except where we replace log price with log sales in some regressions.} If the elasticity of demand is $\sigma$, then $\beta = -\sigma + 1$.

The raw change in fees $\Delta \tau_j$ is mechanically related to changes in consumer composition at the firm level (e.g., if growing firms attract more debit users).
Therefore, we construct an instrument for $\Delta \tau_j$ based on pre-Durbin merchant characteristics and post-Durbin geographic characteristics.
Formally, we construct the instrument as follows:

\begin{align}\widetilde{\Delta\tau_{j}} & = r_{z(j)}\times d_{j}\times\left(\tau^{\text{Reg Debit,Post}}(t_{j})-\tau^{\text{Exempt Debit,Post}}(t_{j})\right) \\
    \tau^{\text{Reg Debit,Post}}(t_{j}) & = 0.0005 + \frac{0.22}{t_j} \\
    \tau^{\text{Exempt Debit,Post}}(t_{j}) & = 0.0119 + \frac{0.10}{t_j}
\end{align}
Given this design, we include $\tau^{\text{Reg Debit,Post}}(t_{j}) - \tau^{\text{Exempt Debit,Post}}(t_{j})$ and $d_j$ as controls in $X_j$.

One potential concern is that the Durbin Amendment may have shifted aggregate payment behavior toward debit cards, since regulated debit became cheaper for merchants.
If so, pre-Durbin payment composition would be less predictive of post-Durbin fee changes, weakening the instrument.
In practice, this is not a major issue: our first-stage F-statistics exceed 20,000, indicating that pre-Durbin characteristics remain highly predictive of fee changes.

We report results for the one-year change in sales in columns (1)-(2) of Table \ref{tab:durbin-iv}.
Column (1) presents OLS estimates, while column (2) presents our IV estimates.
The OLS results indicate that a one-percentage-point increase in average interchange fees leads to approximately a 21.1\% decrease in sales.
However, as discussed previously, there are many reasons why increases in interchange rates are mechanically associated with lower sales.
Column (2) then presents IV estimates that a one-percentage-point increase in average interchange fees causes a 6.9\% decrease in sales.
Columns (3) and (4) show the two-year effects.
Assuming full pass-through of interchange fees to prices, these estimates imply a demand elasticity of around 7.9.
This elasticity is large---higher than the elasticity of around $5$ found in \citet{Sullivan2024}'s study of restaurants using detailed transaction data.
One explanation is that merchants respond to interchange fee changes not only by adjusting posted prices but also by imposing cash discounts or card surcharges that steer customers across payment methods, thereby amplifying the measured sales response.

To test this explanation, we re-estimate our Durbin-Amendment effects using only the ten states that prohibit any form of card surcharging (Table \ref{tab:durbin-iv-nosurcharge}).\footnote{These states are California, Colorado, Connecticut, Florida, Kansas, Massachusetts, Maine, New York, Oklahoma, and Texas.}
If payment-method steering were driving the sales responses, we would expect these effects to dissipate where surcharges are illegal.
The implied elasticity in the no-surcharge subsample is 4.8---smaller than the full-sample estimate, suggesting that surcharging does contribute to the large elasticities.
However, the core result remains: Interchange fee changes cause economically meaningful changes in restaurant sales even where surcharging is prohibited.

This decline in sales could reflect either a price increase or a decline in product quality (e.g., a rise in the quality-adjusted price).
Although we do not observe the prices set by the restaurants, we do observe average transaction sizes as a proxy for price.
Average transaction size is a potential proxy for price if consumers respond to higher restaurant prices by going to a restaurant less often while ordering the same quantities for any given trip.
Column (5) shows that increases in interchange are strong negative predictors of average transaction size.
This reflects the mechanical negative correlation between transaction size and percentage interchange rates that arises when interchange fees have fixed components.
Column (6) shows that our IV estimate instead finds that a one-percentage-point increase in interchange causes a roughly one-for-one increase in prices, but with substantial standard errors.
These results align with the predictions of \citet{Amiti.etal2019}, who show that small firms tend to pass through cost shocks one-for-one to prices.

\begin{figure}[H]
  \centering
  \caption{Explaining the Heterogeneous Impact of the Durbin Amendment Across Restaurants}
  \label{fig:durbin-het-inputs}

  \begin{subfigure}[t]{0.7\textwidth}
    \centering
    \includegraphics[width=\linewidth]{tmp_images/___fiserv_output_graphs_durbin_ticket_size_effect_pdf.png}
    \caption{Variation by Ticket Size}
    \label{fig:durbin-het-ticket}
  \end{subfigure}

  \vspace{0.6em}

  \begin{subfigure}[t]{0.7\textwidth}
    \centering
    \includegraphics[width=\linewidth]{tmp_images/___fiserv_output_graphs_durbin_regulated_share_effect_zip_pdf.png}
    \caption{Variation by Share of Regulated Debit Cards}
    \label{fig:durbin-het-regulated}
  \end{subfigure}

  \par\textit{Notes:} Figure \ref{fig:durbin-het-inputs} illustrates the heterogeneous impact of the Durbin Amendment across merchants in the restaurant sector (MCC codes 5812 and 5814) and across regions.
Panel \ref{fig:durbin-het-ticket} shows a binscatter of pre-Durbin average ticket size against pre- and post-Durbin debit interchange rates.
Ticket size is calculated as the total dollar value of debit card transactions divided by the total count of debit card transactions in the October 2010--September 2011 window preceding the Durbin Amendment; observations are at the merchant-year level.
Panel \ref{fig:durbin-het-regulated} shows a binscatter of the ZIP-code-level post-Durbin share of debit cards that earn regulated interchange against pre- and post-Durbin debit interchange rates.

\end{figure}

\begin{table}[H]
\centering
\caption{Summary Statistics: Restaurant Settlement Panel 2010--2013}
\label{tab:durbin-summary}

{
\def\sym#1{\ifmmode^{#1}\else\(^{#1}\)\fi}
\begin{tabular}{l*{1}{cccc}}
\hline
                    &           N&        Mean&          SD&      Median\\
\hline
Annual Sales (M)    &     268,143&        0.65&        1.48&        0.38\\
Share of Sales on Debit&     268,143&        0.66&        0.19&        0.71\\
Regulated Share of Debit in Zip Code&     171,736&        0.49&        0.20&        0.49\\
Average Ticket Size of Debit Transactions&     268,143&       19.21&       16.71&       12.65\\
\hline
\end{tabular}
}
 
\par\textit{Notes:} Table shows summary statistics at the merchant-year level from our 2010--2013 panel of restaurant sales. Regulated debit statistics are computed for 2011 and onward, for merchant-years with valid ZIP codes and positive debit transaction volume. Interchange rates are calculated as total interchange within a category divided by total sales in that category.
\end{table}

\clearpage

\begin{table}[H]
\caption{The Effects of Changes in Interchange Fees on Sales and Prices\label{tab:durbin-iv}}

\begin{centering}
{\tiny{
\def\sym#1{\ifmmode^{#1}\else\(^{#1}\)\fi}
\begin{tabular}{l*{8}{c}}
\toprule
                    &\multicolumn{2}{c}{1Y Sales}               &\multicolumn{2}{c}{2Y Sales}               &\multicolumn{2}{c}{1Y Price}               &\multicolumn{2}{c}{2Y Price}               \\
                    &\multicolumn{1}{c}{OLS}&\multicolumn{1}{c}{IV}&\multicolumn{1}{c}{OLS}&\multicolumn{1}{c}{IV}&\multicolumn{1}{c}{OLS}&\multicolumn{1}{c}{IV}&\multicolumn{1}{c}{OLS}&\multicolumn{1}{c}{IV}\\
                    &         (1)         &         (1)         &         (1)         &         (1)         &         (1)         &         (1)         &         (1)         &         (1)         \\
\midrule
Change in average ICG&      -21.12***&       -6.93***&      -24.81***&       -5.56***&       -5.03***&        1.15***&       -6.00***&        1.37***\\
                    &      (0.50)         &      (0.96)         &      (0.63)         &      (1.20)         &      (0.09)         &      (0.18)         &      (0.12)         &      (0.24)         \\

Debit share in pre-period&       -0.12***&       -0.09***&       -0.09***&       -0.07***&        0.00         &        0.01***&        0.02***&        0.03***\\
                    &      (0.00)         &      (0.01)         &      (0.01)         &      (0.01)         &      (0.00)         &      (0.00)         &      (0.00)         &      (0.00)         \\

Predicted change in debit fees&       19.34***&       10.72***&       22.32***&       12.77***&        3.85***&        0.10         &        3.79***&        0.13         \\
                    &      (0.36)         &      (0.61)         &      (0.40)         &      (0.65)         &      (0.07)         &      (0.12)         &      (0.08)         &      (0.13)         \\

Constant            &        0.18***&                     &        0.23***&                     &        0.02***&                     &        0.02***&                     \\
                    &      (0.00)         &                     &      (0.00)         &                     &      (0.00)         &                     &      (0.00)         &                     \\

Zip FE              &         Yes         &         Yes         &         Yes         &         Yes         &         Yes         &         Yes         &         Yes         &         Yes         \\
\midrule
N                   &     175,177         &     175,177         &     135,867         &     135,867         &     175,177         &     175,177         &     135,867         &     135,867         \\
1st Stage F Stat    &                     &      62,615         &                     &      48,970         &                     &      62,615         &                     &      48,970         \\
\bottomrule
\multicolumn{9}{l}{ Standard errors in parentheses}\\
\multicolumn{9}{l}{ * \(p<0.05\), ** \(p<0.01\), *** \(p<0.001\)}\\
\end{tabular}
}
 }
\par\end{centering}
\par\textit{Notes: } Table \ref{tab:durbin-iv}
presents instrumental variables estimates of the effects of interchange fee changes on restaurant sales and prices in the one- and two-year windows after the Durbin Amendment. The specification compares merchants in areas with and without large, regulated banks. Key controls for the instrumental variables strategy include (1) the predicted change in debit fees based on the fee schedule and pre-Durbin transaction size and (2) the pre-Durbin debit share. Log ticket size is measured as the log of the ratio of sales to number of transactions and serves as a proxy for prices.
\end{table}

\clearpage

\begin{table}[H]
\caption{The Effects of Changes in Interchange Fees on Sales and Prices\label{tab:durbin-iv-nosurcharge}}

\begin{centering}
{\tiny{
\def\sym#1{\ifmmode^{#1}\else\(^{#1}\)\fi}
\begin{tabular}{l*{8}{c}}
\toprule
                    &\multicolumn{2}{c}{1Y Sales}               &\multicolumn{2}{c}{2Y Sales}               &\multicolumn{2}{c}{1Y Price}               &\multicolumn{2}{c}{2Y Price}               \\
                    &\multicolumn{1}{c}{OLS}&\multicolumn{1}{c}{IV}&\multicolumn{1}{c}{OLS}&\multicolumn{1}{c}{IV}&\multicolumn{1}{c}{OLS}&\multicolumn{1}{c}{IV}&\multicolumn{1}{c}{OLS}&\multicolumn{1}{c}{IV}\\
                    &         (1)         &         (1)         &         (1)         &         (1)         &         (1)         &         (1)         &         (1)         &         (1)         \\
\midrule
Change in average ICG&      -18.93***&       -3.75*  &      -25.79***&       -3.25         &       -5.04***&        1.16***&       -5.85***&        1.68***\\
                    &      (0.77)         &      (1.59)         &      (0.97)         &      (1.98)         &      (0.14)         &      (0.30)         &      (0.19)         &      (0.39)         \\

Debit share in pre-period&       -0.12***&       -0.08***&       -0.10***&       -0.06***&       -0.00         &        0.01***&        0.01***&        0.03***\\
                    &      (0.01)         &      (0.01)         &      (0.01)         &      (0.01)         &      (0.00)         &      (0.00)         &      (0.00)         &      (0.00)         \\

Predicted change in debit fees&       18.30***&        8.92***&       23.21***&       11.81***&        3.89***&        0.06         &        3.73***&       -0.08         \\
                    &      (0.56)         &      (1.03)         &      (0.63)         &      (1.08)         &      (0.11)         &      (0.19)         &      (0.12)         &      (0.21)         \\

Constant            &        0.18***&                     &        0.25***&                     &        0.02***&                     &        0.02***&                     \\
                    &      (0.01)         &                     &      (0.01)         &                     &      (0.00)         &                     &      (0.00)         &                     \\

Zip FE              &         Yes         &         Yes         &         Yes         &         Yes         &         Yes         &         Yes         &         Yes         &         Yes         \\
\midrule
N                   &      72,819         &      72,819         &      55,798         &      55,798         &      72,819         &      72,819         &      55,798         &      55,798         \\
1st Stage F Stat    &                     &      21,108         &                     &      16,593         &                     &      21,108         &                     &      16,593         \\
\bottomrule
\multicolumn{9}{l}{ Standard errors in parentheses}\\
\multicolumn{9}{l}{ * \(p<0.05\), ** \(p<0.01\), *** \(p<0.001\)}\\
\end{tabular}
}
 }
\par\end{centering}
\par\textit{Notes: } Table \ref{tab:durbin-iv-nosurcharge}
presents instrumental variables estimates of the effects of interchange fee changes on restaurant sales and prices in the one- and two-year windows after the Durbin Amendment on a subset of states that prohibit card surcharging.
\end{table}

\end{document}
